% Options for packages loaded elsewhere
\PassOptionsToPackage{unicode}{hyperref}
\PassOptionsToPackage{hyphens}{url}
%
\documentclass[
]{book}
\usepackage{amsmath,amssymb}
\usepackage{iftex}
\ifPDFTeX
  \usepackage[T1]{fontenc}
  \usepackage[utf8]{inputenc}
  \usepackage{textcomp} % provide euro and other symbols
\else % if luatex or xetex
  \usepackage{unicode-math} % this also loads fontspec
  \defaultfontfeatures{Scale=MatchLowercase}
  \defaultfontfeatures[\rmfamily]{Ligatures=TeX,Scale=1}
\fi
\usepackage{lmodern}
\ifPDFTeX\else
  % xetex/luatex font selection
\fi
% Use upquote if available, for straight quotes in verbatim environments
\IfFileExists{upquote.sty}{\usepackage{upquote}}{}
\IfFileExists{microtype.sty}{% use microtype if available
  \usepackage[]{microtype}
  \UseMicrotypeSet[protrusion]{basicmath} % disable protrusion for tt fonts
}{}
\makeatletter
\@ifundefined{KOMAClassName}{% if non-KOMA class
  \IfFileExists{parskip.sty}{%
    \usepackage{parskip}
  }{% else
    \setlength{\parindent}{0pt}
    \setlength{\parskip}{6pt plus 2pt minus 1pt}}
}{% if KOMA class
  \KOMAoptions{parskip=half}}
\makeatother
\usepackage{xcolor}
\usepackage{color}
\usepackage{fancyvrb}
\newcommand{\VerbBar}{|}
\newcommand{\VERB}{\Verb[commandchars=\\\{\}]}
\DefineVerbatimEnvironment{Highlighting}{Verbatim}{commandchars=\\\{\}}
% Add ',fontsize=\small' for more characters per line
\usepackage{framed}
\definecolor{shadecolor}{RGB}{248,248,248}
\newenvironment{Shaded}{\begin{snugshade}}{\end{snugshade}}
\newcommand{\AlertTok}[1]{\textcolor[rgb]{0.94,0.16,0.16}{#1}}
\newcommand{\AnnotationTok}[1]{\textcolor[rgb]{0.56,0.35,0.01}{\textbf{\textit{#1}}}}
\newcommand{\AttributeTok}[1]{\textcolor[rgb]{0.13,0.29,0.53}{#1}}
\newcommand{\BaseNTok}[1]{\textcolor[rgb]{0.00,0.00,0.81}{#1}}
\newcommand{\BuiltInTok}[1]{#1}
\newcommand{\CharTok}[1]{\textcolor[rgb]{0.31,0.60,0.02}{#1}}
\newcommand{\CommentTok}[1]{\textcolor[rgb]{0.56,0.35,0.01}{\textit{#1}}}
\newcommand{\CommentVarTok}[1]{\textcolor[rgb]{0.56,0.35,0.01}{\textbf{\textit{#1}}}}
\newcommand{\ConstantTok}[1]{\textcolor[rgb]{0.56,0.35,0.01}{#1}}
\newcommand{\ControlFlowTok}[1]{\textcolor[rgb]{0.13,0.29,0.53}{\textbf{#1}}}
\newcommand{\DataTypeTok}[1]{\textcolor[rgb]{0.13,0.29,0.53}{#1}}
\newcommand{\DecValTok}[1]{\textcolor[rgb]{0.00,0.00,0.81}{#1}}
\newcommand{\DocumentationTok}[1]{\textcolor[rgb]{0.56,0.35,0.01}{\textbf{\textit{#1}}}}
\newcommand{\ErrorTok}[1]{\textcolor[rgb]{0.64,0.00,0.00}{\textbf{#1}}}
\newcommand{\ExtensionTok}[1]{#1}
\newcommand{\FloatTok}[1]{\textcolor[rgb]{0.00,0.00,0.81}{#1}}
\newcommand{\FunctionTok}[1]{\textcolor[rgb]{0.13,0.29,0.53}{\textbf{#1}}}
\newcommand{\ImportTok}[1]{#1}
\newcommand{\InformationTok}[1]{\textcolor[rgb]{0.56,0.35,0.01}{\textbf{\textit{#1}}}}
\newcommand{\KeywordTok}[1]{\textcolor[rgb]{0.13,0.29,0.53}{\textbf{#1}}}
\newcommand{\NormalTok}[1]{#1}
\newcommand{\OperatorTok}[1]{\textcolor[rgb]{0.81,0.36,0.00}{\textbf{#1}}}
\newcommand{\OtherTok}[1]{\textcolor[rgb]{0.56,0.35,0.01}{#1}}
\newcommand{\PreprocessorTok}[1]{\textcolor[rgb]{0.56,0.35,0.01}{\textit{#1}}}
\newcommand{\RegionMarkerTok}[1]{#1}
\newcommand{\SpecialCharTok}[1]{\textcolor[rgb]{0.81,0.36,0.00}{\textbf{#1}}}
\newcommand{\SpecialStringTok}[1]{\textcolor[rgb]{0.31,0.60,0.02}{#1}}
\newcommand{\StringTok}[1]{\textcolor[rgb]{0.31,0.60,0.02}{#1}}
\newcommand{\VariableTok}[1]{\textcolor[rgb]{0.00,0.00,0.00}{#1}}
\newcommand{\VerbatimStringTok}[1]{\textcolor[rgb]{0.31,0.60,0.02}{#1}}
\newcommand{\WarningTok}[1]{\textcolor[rgb]{0.56,0.35,0.01}{\textbf{\textit{#1}}}}
\usepackage{longtable,booktabs,array}
\usepackage{calc} % for calculating minipage widths
% Correct order of tables after \paragraph or \subparagraph
\usepackage{etoolbox}
\makeatletter
\patchcmd\longtable{\par}{\if@noskipsec\mbox{}\fi\par}{}{}
\makeatother
% Allow footnotes in longtable head/foot
\IfFileExists{footnotehyper.sty}{\usepackage{footnotehyper}}{\usepackage{footnote}}
\makesavenoteenv{longtable}
\usepackage{graphicx}
\makeatletter
\newsavebox\pandoc@box
\newcommand*\pandocbounded[1]{% scales image to fit in text height/width
  \sbox\pandoc@box{#1}%
  \Gscale@div\@tempa{\textheight}{\dimexpr\ht\pandoc@box+\dp\pandoc@box\relax}%
  \Gscale@div\@tempb{\linewidth}{\wd\pandoc@box}%
  \ifdim\@tempb\p@<\@tempa\p@\let\@tempa\@tempb\fi% select the smaller of both
  \ifdim\@tempa\p@<\p@\scalebox{\@tempa}{\usebox\pandoc@box}%
  \else\usebox{\pandoc@box}%
  \fi%
}
% Set default figure placement to htbp
\def\fps@figure{htbp}
\makeatother
\setlength{\emergencystretch}{3em} % prevent overfull lines
\providecommand{\tightlist}{%
  \setlength{\itemsep}{0pt}\setlength{\parskip}{0pt}}
\setcounter{secnumdepth}{5}
\usepackage{booktabs}
\usepackage{amsthm}
\makeatletter
\def\thm@space@setup{%
  \thm@preskip=8pt plus 2pt minus 4pt
  \thm@postskip=\thm@preskip
}
\makeatother
\usepackage[]{natbib}
\bibliographystyle{apalike}
\usepackage{bookmark}
\IfFileExists{xurl.sty}{\usepackage{xurl}}{} % add URL line breaks if available
\urlstyle{same}
\hypersetup{
  pdftitle={Research in Computational Biology book},
  pdfauthor={Gepoliano Chaves},
  hidelinks,
  pdfcreator={LaTeX via pandoc}}

\title{Research in Computational Biology book}
\author{Gepoliano Chaves}
\date{2026-07-20}

\begin{document}
\maketitle

{
\setcounter{tocdepth}{1}
\tableofcontents
}
\chapter{Prerequisites}\label{prerequisites}

This is a \emph{sample} book written in \textbf{Markdown}. You can use anything that Pandoc's Markdown supports, e.g., a math equation \(a^2 + b^2 = c^2\).

The \textbf{bookdown} package can be installed from CRAN or Github:

\begin{Shaded}
\begin{Highlighting}[]
\FunctionTok{install.packages}\NormalTok{(}\StringTok{"bookdown"}\NormalTok{)}
\CommentTok{\# or the development version}
\CommentTok{\# devtools::install\_github("rstudio/bookdown")}
\end{Highlighting}
\end{Shaded}

Remember each Rmd file contains one and only one chapter, and a chapter is defined by the first-level heading \texttt{\#}.

To compile this example to PDF, you need XeLaTeX. You are recommended to install TinyTeX (which includes XeLaTeX): \url{https://yihui.org/tinytex/}.

\chapter{Basics of R}\label{basics-of-r}

\section{Basic Steps}\label{basic-steps}

\begin{itemize}
\item
  R can be accessed by clicking an icon or entering the command ``R'' at the system command line, also refered to as Terminal
\item
  This produces a console window or causes R to start up as an interactive program at the current terminal window
\item
  R works fundamentally by a question-and-answer model: enter a line with a command and press Enter
\item
  Then the program does something relevant
\item
  In this course, we will use RStudio and a file format called R Markdown, which allows us to code using not only R, but bash and Python, two other languages we may refer to as the course progresses
\end{itemize}

\section{Library ISwR}\label{library-iswr}

\begin{itemize}
\item
  This notebook used the book Introductory Statistics with R, by Peter Dalgaard as a reference
\item
  For the book, R library ISwR (Introductory Statistics with R) can be freely downloaded
\item
  All examples in the book used as reference should run provided that ISwR library is not only installed by also loaded into the current search path
\item
  Library can be installed and loaded by typing the following command into an R chunk
\end{itemize}

\begin{Shaded}
\begin{Highlighting}[]
\FunctionTok{install.packages}\NormalTok{(}\StringTok{"ISwR"}\NormalTok{)}
\FunctionTok{library}\NormalTok{(}\StringTok{"ISwR"}\NormalTok{)}
\end{Highlighting}
\end{Shaded}

\section{Plot Random Numbers}\label{plot-random-numbers}

\begin{itemize}
\item
  For a first impression of what R can do, let´s try plotting a graph
\item
  Need to insert the R chunk and use the plot function
\end{itemize}

\begin{Shaded}
\begin{Highlighting}[]
\FunctionTok{plot}\NormalTok{(}\FunctionTok{rnorm}\NormalTok{(}\DecValTok{1000}\NormalTok{))}
\end{Highlighting}
\end{Shaded}

\pandocbounded{\includegraphics[keepaspectratio]{bookdown-demo_files/figure-latex/unnamed-chunk-3-1.pdf}}

\section{A Potent Calculator}\label{a-potent-calculator}

\begin{itemize}
\tightlist
\item
  R can process simple and complex arithmetic expression and produce a result for the user
\end{itemize}

\begin{Shaded}
\begin{Highlighting}[]
\DecValTok{2} \SpecialCharTok{+} \DecValTok{2}
\end{Highlighting}
\end{Shaded}

\begin{verbatim}
## [1] 4
\end{verbatim}

\begin{itemize}
\tightlist
\item
  R can also be used to do other standard calculations. Here is how to calculate e to the power of -2
\end{itemize}

\begin{Shaded}
\begin{Highlighting}[]
\FunctionTok{exp}\NormalTok{(}\SpecialCharTok{{-}}\DecValTok{2}\NormalTok{)}
\end{Highlighting}
\end{Shaded}

\begin{verbatim}
## [1] 0.1353353
\end{verbatim}

\begin{itemize}
\item
  Other than the R chunks, these calculations can be made using the RStudio Console
\item
  In RStusio, when using the R Markdown format (.Rmd), we can insert R, Bash (terminal) and Python chunks
\item
  Proceed to illustrate that in the RStudio IDE
\end{itemize}

\section{Assignments}\label{assignments}

\begin{itemize}
\item
  Assignments are made based on the necessity to store the results of calculations and use these results in downstream processing steps in an entire algorithm or pipeline
\item
  Like other languages, R has symbolic \textbf{variables}: names that can be used to represent values
\item
  To assign the value 2 to variable x, we can enter the following command
\end{itemize}

\begin{Shaded}
\begin{Highlighting}[]
\NormalTok{x }\OtherTok{\textless{}{-}} \DecValTok{2}
\end{Highlighting}
\end{Shaded}

\begin{itemize}
\item
  The character \textless- is called the assignment operator
\item
  There is no immediate visible result, but from now on, x has the value 2 and can be used in subsequent arithmetic operations
\item
  We can ``ask'' R which value is stores in the x variable
\end{itemize}

\begin{Shaded}
\begin{Highlighting}[]
\NormalTok{x}
\end{Highlighting}
\end{Shaded}

\begin{verbatim}
## [1] 2
\end{verbatim}

\section{Operations after assignment of variable}\label{operations-after-assignment-of-variable}

\begin{itemize}
\tightlist
\item
  Below, our variable x, is used to perform other calculations
\end{itemize}

\begin{Shaded}
\begin{Highlighting}[]
\NormalTok{x}
\end{Highlighting}
\end{Shaded}

\begin{verbatim}
## [1] 2
\end{verbatim}

\begin{itemize}
\tightlist
\item
  Addition
\end{itemize}

\begin{Shaded}
\begin{Highlighting}[]
\NormalTok{x }\SpecialCharTok{+}\NormalTok{ x}
\end{Highlighting}
\end{Shaded}

\begin{verbatim}
## [1] 4
\end{verbatim}

\begin{itemize}
\tightlist
\item
  Multiplication
\end{itemize}

\begin{Shaded}
\begin{Highlighting}[]
\DecValTok{5}\SpecialCharTok{*}\NormalTok{x}
\end{Highlighting}
\end{Shaded}

\begin{verbatim}
## [1] 10
\end{verbatim}

\begin{itemize}
\tightlist
\item
  Potentiation
\end{itemize}

\begin{Shaded}
\begin{Highlighting}[]
\NormalTok{x}\SpecialCharTok{\^{}}\DecValTok{3} 
\end{Highlighting}
\end{Shaded}

\begin{verbatim}
## [1] 8
\end{verbatim}

\section{Vectorized Arithmetic}\label{vectorized-arithmetic}

\begin{itemize}
\item
  It is not useful to use one number at a time to run statistics
\item
  One strenght of R is that it can handle entire data vectors as single objects
\item
  A data vector is an array of numbers and a vector variable can be constructed like this
\end{itemize}

\begin{Shaded}
\begin{Highlighting}[]
\NormalTok{weight }\OtherTok{\textless{}{-}} \FunctionTok{c}\NormalTok{(}\DecValTok{60}\NormalTok{, }\DecValTok{72}\NormalTok{, }\DecValTok{57}\NormalTok{, }\DecValTok{90}\NormalTok{, }\DecValTok{95}\NormalTok{, }\DecValTok{72}\NormalTok{)}

\DocumentationTok{\#\# To look at the vector variable, just type its name again}
\NormalTok{weight}
\end{Highlighting}
\end{Shaded}

\begin{verbatim}
## [1] 60 72 57 90 95 72
\end{verbatim}

\begin{itemize}
\item
  You can do calculations with vectors just like ordinary numbers, so long as they have the same length
\item
  One exception to this rule that we will see will be when we use the mean of weigths of persons (represented by xbar)
\item
  In that case, the mean will be one single number, which will be subtracted from each sample value
\item
  Suppose the weight vector indicates the weight of men in kilograms
\item
  One simple formula to indicate whether a person is obese or not, is the body mass index (BMI)
\item
  BMI is calculated by dividing the person´s weight by the square of their height, in meters
\item
  Therefore, in R, we need to have a vector with the height values to calculate the bmi vector, containing the body-mass index for the individuals indicated in the weight vector
\end{itemize}

\begin{Shaded}
\begin{Highlighting}[]
\NormalTok{height }\OtherTok{\textless{}{-}} \FunctionTok{c}\NormalTok{(}\FloatTok{1.75}\NormalTok{, }\FloatTok{1.80}\NormalTok{, }\FloatTok{1.65}\NormalTok{, }\FloatTok{1.90}\NormalTok{, }\FloatTok{1.74}\NormalTok{, }\FloatTok{1.91}\NormalTok{)}
\NormalTok{bmi }\OtherTok{\textless{}{-}}\NormalTok{ weight}\SpecialCharTok{/}\NormalTok{height}\SpecialCharTok{\^{}}\DecValTok{2}
\NormalTok{bmi}
\end{Highlighting}
\end{Shaded}

\begin{verbatim}
## [1] 19.59184 22.22222 20.93664 24.93075 31.37799 19.73630
\end{verbatim}

\section{Calculate the Mean of the variable weight}\label{calculate-the-mean-of-the-variable-weight}

\begin{itemize}
\tightlist
\item
  The mean is calculated by the sum of the observations divided by the total number of observations
\end{itemize}

\(\overline{x}\) = \(\sum x_i\) /n

\begin{itemize}
\tightlist
\item
  Let´s calculate the mean of the variable weight
\end{itemize}

\begin{Shaded}
\begin{Highlighting}[]
\FunctionTok{sum}\NormalTok{(weight)}
\end{Highlighting}
\end{Shaded}

\begin{verbatim}
## [1] 446
\end{verbatim}

\begin{Shaded}
\begin{Highlighting}[]
\FunctionTok{sum}\NormalTok{(weight)}\SpecialCharTok{/}\FunctionTok{length}\NormalTok{(weight)}
\end{Highlighting}
\end{Shaded}

\begin{verbatim}
## [1] 74.33333
\end{verbatim}

\section{Calculate the Standard-Deviation of the variable weight}\label{calculate-the-standard-deviation-of-the-variable-weight}

\begin{itemize}
\tightlist
\item
  Standard-deviation can be calculated with the following equation
\end{itemize}

\[ SD = \sqrt{(\sum (x_i - \overline{x})^2)/(n-1)}\]

\begin{itemize}
\tightlist
\item
  xbar, the mean of variable weight, can be calculated using the sum and length of variable weight
\end{itemize}

\begin{Shaded}
\begin{Highlighting}[]
\NormalTok{xbar }\OtherTok{\textless{}{-}} \FunctionTok{sum}\NormalTok{(weight)}\SpecialCharTok{/}\FunctionTok{length}\NormalTok{(weight)}
\NormalTok{xbar}
\end{Highlighting}
\end{Shaded}

\begin{verbatim}
## [1] 74.33333
\end{verbatim}

\begin{itemize}
\tightlist
\item
  Now we can calculate the difference of each replicate in the weight variable and the mean of the weight variable, one by one
\end{itemize}

\begin{Shaded}
\begin{Highlighting}[]
\NormalTok{weight }\SpecialCharTok{{-}}\NormalTok{ xbar}
\end{Highlighting}
\end{Shaded}

\begin{verbatim}
## [1] -14.333333  -2.333333 -17.333333  15.666667  20.666667  -2.333333
\end{verbatim}

\begin{itemize}
\item
  Notice how R uses xbar, which has length one, to calculate the new weight - xbar data vector
\item
  xbar is recycled and subracted from every element in the weight variable (weight data vector)
\end{itemize}

\section{Calculation of the Standard Deviation}\label{calculation-of-the-standard-deviation}

\begin{itemize}
\tightlist
\item
  Calculate the squared deviations
\end{itemize}

\begin{Shaded}
\begin{Highlighting}[]
\NormalTok{(weight }\SpecialCharTok{{-}}\NormalTok{ xbar)}\SpecialCharTok{\^{}}\DecValTok{2}
\end{Highlighting}
\end{Shaded}

\begin{verbatim}
## [1] 205.444444   5.444444 300.444444 245.444444 427.111111   5.444444
\end{verbatim}

\begin{itemize}
\tightlist
\item
  Calculate the sum of squared deviations
\end{itemize}

\begin{Shaded}
\begin{Highlighting}[]
\FunctionTok{sum}\NormalTok{((weight }\SpecialCharTok{{-}}\NormalTok{ xbar)}\SpecialCharTok{\^{}}\DecValTok{2}\NormalTok{)}
\end{Highlighting}
\end{Shaded}

\begin{verbatim}
## [1] 1189.333
\end{verbatim}

Calculate the standard deviation

\begin{Shaded}
\begin{Highlighting}[]
\FunctionTok{sqrt}\NormalTok{(}
  \FunctionTok{sum}\NormalTok{(}
\NormalTok{    (weight }\SpecialCharTok{{-}}\NormalTok{ xbar)}\SpecialCharTok{\^{}}\DecValTok{2}\SpecialCharTok{/}
\NormalTok{      (}\FunctionTok{length}\NormalTok{(weight)}\SpecialCharTok{{-}}\DecValTok{1}\NormalTok{)}
\NormalTok{)}
\NormalTok{)}
\end{Highlighting}
\end{Shaded}

\begin{verbatim}
## [1] 15.42293
\end{verbatim}

\section{Standard Statistical Procedures}\label{standard-statistical-procedures}

\begin{itemize}
\item
  It is a standard medical practice to access whether a person is obese or not using validated scientific criteria
\item
  As a simple procedure to show this concept, let's assume that an individual with a normal weight should have a BMI in the range 20-25
\item
  We want to know if our data deviates from the normal range of BMI
\item
  In R, this can be done using a statistical test called t-test
\item
  You do not need to understand what a t-test is, just remember that is is used to evaluate the distribution of sample values compared to the normal distribution
\item
  You can use a one-sample t-test to assess whether the six persons' BMI can be assumed to have mean 22.5 given that they come from a normal distribution
\item
  You can do that using the function t.test
\end{itemize}

\section{Standard Statistical Procedures: t-test}\label{standard-statistical-procedures-t-test}

\begin{Shaded}
\begin{Highlighting}[]
\FunctionTok{t.test}\NormalTok{(}
\NormalTok{  bmi,}
  \AttributeTok{mu =} \FloatTok{22.5}
\NormalTok{)}
\end{Highlighting}
\end{Shaded}

\begin{verbatim}
## 
##  One Sample t-test
## 
## data:  bmi
## t = 0.34488, df = 5, p-value = 0.7442
## alternative hypothesis: true mean is not equal to 22.5
## 95 percent confidence interval:
##  18.41734 27.84791
## sample estimates:
## mean of x 
##  23.13262
\end{verbatim}

\section{Plot Graphics}\label{plot-graphics}

\begin{itemize}
\tightlist
\item
  Let's now plot a scatterplot of the height and weight of individuals
\end{itemize}

\begin{Shaded}
\begin{Highlighting}[]
\FunctionTok{plot}\NormalTok{(height, weight)}
\end{Highlighting}
\end{Shaded}

\pandocbounded{\includegraphics[keepaspectratio]{bookdown-demo_files/figure-latex/unnamed-chunk-19-1.pdf}}

\section{Plot Graphics Modifying the plotting character}\label{plot-graphics-modifying-the-plotting-character}

\begin{itemize}
\item
  You will frequently want to modify drawing of your graphs in various ways
\item
  One way is usig the parameter ``plotting character'', pch
\end{itemize}

\begin{Shaded}
\begin{Highlighting}[]
\FunctionTok{plot}\NormalTok{(height, weight, }\AttributeTok{pch =}\DecValTok{2}\NormalTok{)}
\end{Highlighting}
\end{Shaded}

\pandocbounded{\includegraphics[keepaspectratio]{bookdown-demo_files/figure-latex/unnamed-chunk-20-1.pdf}}

\section{Plot an expected Line for weights at BMI of 22.5}\label{plot-an-expected-line-for-weights-at-bmi-of-22.5}

\begin{itemize}
\tightlist
\item
  The term for weight at BMI of 22.5 is:
\end{itemize}

\begin{Shaded}
\begin{Highlighting}[]
\FloatTok{22.5}\SpecialCharTok{*}\NormalTok{(height)}\SpecialCharTok{\^{}}\DecValTok{2}
\end{Highlighting}
\end{Shaded}

\begin{itemize}
\item
  This term can be used to compute the line of expected weights at a BMI of 22.5 using the lines() function
\item
  As a note, we can't use the lines() function without first creating a plot in R.
\end{itemize}

\begin{Shaded}
\begin{Highlighting}[]
\FunctionTok{plot}\NormalTok{(height, weight, }\AttributeTok{pch =}\DecValTok{2}\NormalTok{)}
\NormalTok{hh }\OtherTok{\textless{}{-}} \FunctionTok{c}\NormalTok{(}\FloatTok{1.65}\NormalTok{, }\FloatTok{1.70}\NormalTok{, }\FloatTok{1.75}\NormalTok{, }\FloatTok{1.80}\NormalTok{, }\FloatTok{1.85}\NormalTok{, }\FloatTok{1.90}\NormalTok{)}
\FunctionTok{lines}\NormalTok{(hh, }\FloatTok{22.5} \SpecialCharTok{*}\NormalTok{ (hh)}\SpecialCharTok{\^{}}\DecValTok{2}\NormalTok{)}
\end{Highlighting}
\end{Shaded}

\pandocbounded{\includegraphics[keepaspectratio]{bookdown-demo_files/figure-latex/unnamed-chunk-22-1.pdf}}

\section{Vectors}\label{vectors}

\begin{itemize}
\item
  The weight and height vectors are called numeric vectors
\item
  Besides numeric vectors, there are numeric and character vectors
\end{itemize}

\section{Using other libraries}\label{using-other-libraries}

\subsection{The iris data-frame}\label{the-iris-data-frame}

\begin{itemize}
\item
  Now we have a better understanding of what R can give us, let us use another library more commonly used datasets
\item
  At times, it is possible that you will need to figure out different ways to install a library to use it
\end{itemize}

\begin{Shaded}
\begin{Highlighting}[]
\FunctionTok{library}\NormalTok{(datasets)}
\end{Highlighting}
\end{Shaded}

\begin{itemize}
\tightlist
\item
  In the next chunk, we access the iris data, and look at a summary of the dataset
\end{itemize}

\begin{Shaded}
\begin{Highlighting}[]
\FunctionTok{data}\NormalTok{(iris)}
\FunctionTok{summary}\NormalTok{(iris)}
\end{Highlighting}
\end{Shaded}

\begin{verbatim}
##   Sepal.Length    Sepal.Width     Petal.Length    Petal.Width   
##  Min.   :4.300   Min.   :2.000   Min.   :1.000   Min.   :0.100  
##  1st Qu.:5.100   1st Qu.:2.800   1st Qu.:1.600   1st Qu.:0.300  
##  Median :5.800   Median :3.000   Median :4.350   Median :1.300  
##  Mean   :5.843   Mean   :3.057   Mean   :3.758   Mean   :1.199  
##  3rd Qu.:6.400   3rd Qu.:3.300   3rd Qu.:5.100   3rd Qu.:1.800  
##  Max.   :7.900   Max.   :4.400   Max.   :6.900   Max.   :2.500  
##        Species  
##  setosa    :50  
##  versicolor:50  
##  virginica :50  
##                 
##                 
## 
\end{verbatim}

\begin{itemize}
\tightlist
\item
  Another form to look at the iris data-frame is typing its name
\end{itemize}

\begin{Shaded}
\begin{Highlighting}[]
\NormalTok{iris}
\end{Highlighting}
\end{Shaded}

\begin{verbatim}
##     Sepal.Length Sepal.Width Petal.Length Petal.Width    Species
## 1            5.1         3.5          1.4         0.2     setosa
## 2            4.9         3.0          1.4         0.2     setosa
## 3            4.7         3.2          1.3         0.2     setosa
## 4            4.6         3.1          1.5         0.2     setosa
## 5            5.0         3.6          1.4         0.2     setosa
## 6            5.4         3.9          1.7         0.4     setosa
## 7            4.6         3.4          1.4         0.3     setosa
## 8            5.0         3.4          1.5         0.2     setosa
## 9            4.4         2.9          1.4         0.2     setosa
## 10           4.9         3.1          1.5         0.1     setosa
## 11           5.4         3.7          1.5         0.2     setosa
## 12           4.8         3.4          1.6         0.2     setosa
## 13           4.8         3.0          1.4         0.1     setosa
## 14           4.3         3.0          1.1         0.1     setosa
## 15           5.8         4.0          1.2         0.2     setosa
## 16           5.7         4.4          1.5         0.4     setosa
## 17           5.4         3.9          1.3         0.4     setosa
## 18           5.1         3.5          1.4         0.3     setosa
## 19           5.7         3.8          1.7         0.3     setosa
## 20           5.1         3.8          1.5         0.3     setosa
## 21           5.4         3.4          1.7         0.2     setosa
## 22           5.1         3.7          1.5         0.4     setosa
## 23           4.6         3.6          1.0         0.2     setosa
## 24           5.1         3.3          1.7         0.5     setosa
## 25           4.8         3.4          1.9         0.2     setosa
## 26           5.0         3.0          1.6         0.2     setosa
## 27           5.0         3.4          1.6         0.4     setosa
## 28           5.2         3.5          1.5         0.2     setosa
## 29           5.2         3.4          1.4         0.2     setosa
## 30           4.7         3.2          1.6         0.2     setosa
## 31           4.8         3.1          1.6         0.2     setosa
## 32           5.4         3.4          1.5         0.4     setosa
## 33           5.2         4.1          1.5         0.1     setosa
## 34           5.5         4.2          1.4         0.2     setosa
## 35           4.9         3.1          1.5         0.2     setosa
## 36           5.0         3.2          1.2         0.2     setosa
## 37           5.5         3.5          1.3         0.2     setosa
## 38           4.9         3.6          1.4         0.1     setosa
## 39           4.4         3.0          1.3         0.2     setosa
## 40           5.1         3.4          1.5         0.2     setosa
## 41           5.0         3.5          1.3         0.3     setosa
## 42           4.5         2.3          1.3         0.3     setosa
## 43           4.4         3.2          1.3         0.2     setosa
## 44           5.0         3.5          1.6         0.6     setosa
## 45           5.1         3.8          1.9         0.4     setosa
## 46           4.8         3.0          1.4         0.3     setosa
## 47           5.1         3.8          1.6         0.2     setosa
## 48           4.6         3.2          1.4         0.2     setosa
## 49           5.3         3.7          1.5         0.2     setosa
## 50           5.0         3.3          1.4         0.2     setosa
## 51           7.0         3.2          4.7         1.4 versicolor
## 52           6.4         3.2          4.5         1.5 versicolor
## 53           6.9         3.1          4.9         1.5 versicolor
## 54           5.5         2.3          4.0         1.3 versicolor
## 55           6.5         2.8          4.6         1.5 versicolor
## 56           5.7         2.8          4.5         1.3 versicolor
## 57           6.3         3.3          4.7         1.6 versicolor
## 58           4.9         2.4          3.3         1.0 versicolor
## 59           6.6         2.9          4.6         1.3 versicolor
## 60           5.2         2.7          3.9         1.4 versicolor
## 61           5.0         2.0          3.5         1.0 versicolor
## 62           5.9         3.0          4.2         1.5 versicolor
## 63           6.0         2.2          4.0         1.0 versicolor
## 64           6.1         2.9          4.7         1.4 versicolor
## 65           5.6         2.9          3.6         1.3 versicolor
## 66           6.7         3.1          4.4         1.4 versicolor
## 67           5.6         3.0          4.5         1.5 versicolor
## 68           5.8         2.7          4.1         1.0 versicolor
## 69           6.2         2.2          4.5         1.5 versicolor
## 70           5.6         2.5          3.9         1.1 versicolor
## 71           5.9         3.2          4.8         1.8 versicolor
## 72           6.1         2.8          4.0         1.3 versicolor
## 73           6.3         2.5          4.9         1.5 versicolor
## 74           6.1         2.8          4.7         1.2 versicolor
## 75           6.4         2.9          4.3         1.3 versicolor
## 76           6.6         3.0          4.4         1.4 versicolor
## 77           6.8         2.8          4.8         1.4 versicolor
## 78           6.7         3.0          5.0         1.7 versicolor
## 79           6.0         2.9          4.5         1.5 versicolor
## 80           5.7         2.6          3.5         1.0 versicolor
## 81           5.5         2.4          3.8         1.1 versicolor
## 82           5.5         2.4          3.7         1.0 versicolor
## 83           5.8         2.7          3.9         1.2 versicolor
## 84           6.0         2.7          5.1         1.6 versicolor
## 85           5.4         3.0          4.5         1.5 versicolor
## 86           6.0         3.4          4.5         1.6 versicolor
## 87           6.7         3.1          4.7         1.5 versicolor
## 88           6.3         2.3          4.4         1.3 versicolor
## 89           5.6         3.0          4.1         1.3 versicolor
## 90           5.5         2.5          4.0         1.3 versicolor
## 91           5.5         2.6          4.4         1.2 versicolor
## 92           6.1         3.0          4.6         1.4 versicolor
## 93           5.8         2.6          4.0         1.2 versicolor
## 94           5.0         2.3          3.3         1.0 versicolor
## 95           5.6         2.7          4.2         1.3 versicolor
## 96           5.7         3.0          4.2         1.2 versicolor
## 97           5.7         2.9          4.2         1.3 versicolor
## 98           6.2         2.9          4.3         1.3 versicolor
## 99           5.1         2.5          3.0         1.1 versicolor
## 100          5.7         2.8          4.1         1.3 versicolor
## 101          6.3         3.3          6.0         2.5  virginica
## 102          5.8         2.7          5.1         1.9  virginica
## 103          7.1         3.0          5.9         2.1  virginica
## 104          6.3         2.9          5.6         1.8  virginica
## 105          6.5         3.0          5.8         2.2  virginica
## 106          7.6         3.0          6.6         2.1  virginica
## 107          4.9         2.5          4.5         1.7  virginica
## 108          7.3         2.9          6.3         1.8  virginica
## 109          6.7         2.5          5.8         1.8  virginica
## 110          7.2         3.6          6.1         2.5  virginica
## 111          6.5         3.2          5.1         2.0  virginica
## 112          6.4         2.7          5.3         1.9  virginica
## 113          6.8         3.0          5.5         2.1  virginica
## 114          5.7         2.5          5.0         2.0  virginica
## 115          5.8         2.8          5.1         2.4  virginica
## 116          6.4         3.2          5.3         2.3  virginica
## 117          6.5         3.0          5.5         1.8  virginica
## 118          7.7         3.8          6.7         2.2  virginica
## 119          7.7         2.6          6.9         2.3  virginica
## 120          6.0         2.2          5.0         1.5  virginica
## 121          6.9         3.2          5.7         2.3  virginica
## 122          5.6         2.8          4.9         2.0  virginica
## 123          7.7         2.8          6.7         2.0  virginica
## 124          6.3         2.7          4.9         1.8  virginica
## 125          6.7         3.3          5.7         2.1  virginica
## 126          7.2         3.2          6.0         1.8  virginica
## 127          6.2         2.8          4.8         1.8  virginica
## 128          6.1         3.0          4.9         1.8  virginica
## 129          6.4         2.8          5.6         2.1  virginica
## 130          7.2         3.0          5.8         1.6  virginica
## 131          7.4         2.8          6.1         1.9  virginica
## 132          7.9         3.8          6.4         2.0  virginica
## 133          6.4         2.8          5.6         2.2  virginica
## 134          6.3         2.8          5.1         1.5  virginica
## 135          6.1         2.6          5.6         1.4  virginica
## 136          7.7         3.0          6.1         2.3  virginica
## 137          6.3         3.4          5.6         2.4  virginica
## 138          6.4         3.1          5.5         1.8  virginica
## 139          6.0         3.0          4.8         1.8  virginica
## 140          6.9         3.1          5.4         2.1  virginica
## 141          6.7         3.1          5.6         2.4  virginica
## 142          6.9         3.1          5.1         2.3  virginica
## 143          5.8         2.7          5.1         1.9  virginica
## 144          6.8         3.2          5.9         2.3  virginica
## 145          6.7         3.3          5.7         2.5  virginica
## 146          6.7         3.0          5.2         2.3  virginica
## 147          6.3         2.5          5.0         1.9  virginica
## 148          6.5         3.0          5.2         2.0  virginica
## 149          6.2         3.4          5.4         2.3  virginica
## 150          5.9         3.0          5.1         1.8  virginica
\end{verbatim}

\subsection{More Visualization of the iris dataset}\label{more-visualization-of-the-iris-dataset}

These are basic R functions that allow us to get general information about a dataframe:

\begin{itemize}
\tightlist
\item
  dim()
\item
  head()
\item
  View()
\item
  class()
\item
  str()
\end{itemize}

\subsection{Basic R functions}\label{basic-r-functions}

\begin{Shaded}
\begin{Highlighting}[]
\FunctionTok{dim}\NormalTok{(iris)}
\FunctionTok{head}\NormalTok{(iris)}
\FunctionTok{View}\NormalTok{(iris)}
\FunctionTok{class}\NormalTok{(iris)}
\FunctionTok{str}\NormalTok{(iris)}
\end{Highlighting}
\end{Shaded}

\begin{itemize}
\tightlist
\item
  Select the number of lines to visualize with the head command
\end{itemize}

\begin{Shaded}
\begin{Highlighting}[]
\FunctionTok{head}\NormalTok{(iris, }\AttributeTok{n =} \DecValTok{10}\NormalTok{)}
\end{Highlighting}
\end{Shaded}

\begin{verbatim}
##    Sepal.Length Sepal.Width Petal.Length Petal.Width Species
## 1           5.1         3.5          1.4         0.2  setosa
## 2           4.9         3.0          1.4         0.2  setosa
## 3           4.7         3.2          1.3         0.2  setosa
## 4           4.6         3.1          1.5         0.2  setosa
## 5           5.0         3.6          1.4         0.2  setosa
## 6           5.4         3.9          1.7         0.4  setosa
## 7           4.6         3.4          1.4         0.3  setosa
## 8           5.0         3.4          1.5         0.2  setosa
## 9           4.4         2.9          1.4         0.2  setosa
## 10          4.9         3.1          1.5         0.1  setosa
\end{verbatim}

\begin{itemize}
\tightlist
\item
  Determine the number of columns and rows in the dataframe:
\end{itemize}

\begin{Shaded}
\begin{Highlighting}[]
\FunctionTok{dim}\NormalTok{(iris)}
\end{Highlighting}
\end{Shaded}

\begin{verbatim}
## [1] 150   5
\end{verbatim}

\begin{itemize}
\tightlist
\item
  Determine the class() that the dataset belongs to
\end{itemize}

\begin{Shaded}
\begin{Highlighting}[]
\FunctionTok{class}\NormalTok{(iris)}
\end{Highlighting}
\end{Shaded}

\begin{verbatim}
## [1] "data.frame"
\end{verbatim}

\section{More Visualization}\label{more-visualization}

\subsection{Scatterplot (includes warnings)}\label{scatterplot-includes-warnings}

\begin{Shaded}
\begin{Highlighting}[]
\FunctionTok{plot}\NormalTok{(}\AttributeTok{data=}\NormalTok{iris, iris}\SpecialCharTok{$}\NormalTok{Sepal.Length, iris}\SpecialCharTok{$}\NormalTok{Sepal.Width) }\DocumentationTok{\#\# R will complain about this command}
\end{Highlighting}
\end{Shaded}

\begin{verbatim}
## Warning in plot.window(...): "data" is not a graphical parameter
\end{verbatim}

\begin{verbatim}
## Warning in plot.xy(xy, type, ...): "data" is not a graphical parameter
\end{verbatim}

\begin{verbatim}
## Warning in axis(side = side, at = at, labels = labels, ...): "data" is not a
## graphical parameter
## Warning in axis(side = side, at = at, labels = labels, ...): "data" is not a
## graphical parameter
\end{verbatim}

\begin{verbatim}
## Warning in box(...): "data" is not a graphical parameter
\end{verbatim}

\begin{verbatim}
## Warning in title(...): "data" is not a graphical parameter
\end{verbatim}

\pandocbounded{\includegraphics[keepaspectratio]{bookdown-demo_files/figure-latex/unnamed-chunk-30-1.pdf}}

\begin{Shaded}
\begin{Highlighting}[]
\FunctionTok{plot}\NormalTok{(iris}\SpecialCharTok{$}\NormalTok{Sepal.Length, iris}\SpecialCharTok{$}\NormalTok{Sepal.Width) }\DocumentationTok{\#\# According to the error, this was plotted using a different synthax}
\end{Highlighting}
\end{Shaded}

\pandocbounded{\includegraphics[keepaspectratio]{bookdown-demo_files/figure-latex/unnamed-chunk-30-2.pdf}}

\subsection{Scatterplot}\label{scatterplot}

Plot without warning and message

\begin{Shaded}
\begin{Highlighting}[]
\FunctionTok{plot}\NormalTok{(}\AttributeTok{data=}\NormalTok{iris, iris}\SpecialCharTok{$}\NormalTok{Sepal.Length, iris}\SpecialCharTok{$}\NormalTok{Sepal.Width) }\DocumentationTok{\#\# R will complain about this command}
\end{Highlighting}
\end{Shaded}

\pandocbounded{\includegraphics[keepaspectratio]{bookdown-demo_files/figure-latex/unnamed-chunk-31-1.pdf}}

\begin{Shaded}
\begin{Highlighting}[]
\FunctionTok{plot}\NormalTok{(iris}\SpecialCharTok{$}\NormalTok{Sepal.Length, iris}\SpecialCharTok{$}\NormalTok{Sepal.Width) }\DocumentationTok{\#\# According to the error, this was plotted using a different synthax}
\end{Highlighting}
\end{Shaded}

\pandocbounded{\includegraphics[keepaspectratio]{bookdown-demo_files/figure-latex/unnamed-chunk-31-2.pdf}}

\subsection{Boxplot}\label{boxplot}

\begin{Shaded}
\begin{Highlighting}[]
\FunctionTok{boxplot}\NormalTok{(}\AttributeTok{data=}\NormalTok{iris, iris}\SpecialCharTok{$}\NormalTok{Sepal.Length, iris}\SpecialCharTok{$}\NormalTok{Sepal.Width)}
\end{Highlighting}
\end{Shaded}

\pandocbounded{\includegraphics[keepaspectratio]{bookdown-demo_files/figure-latex/3. Boxplot-1.pdf}}

\section{Data Visualization with specialized libraries}\label{data-visualization-with-specialized-libraries}

\begin{itemize}
\item
  In R, there are packages designed with the purpose of making good-looking graphics. This is the case with the ggplot2. The chunk below installs the ggplot2 library, loads the library into the R environment and then plots the data present in the iris data-frame.
\item
  You can uncomment the installation line if you need to install it
\end{itemize}

\subsection{Ggplot2}\label{ggplot2}

\begin{Shaded}
\begin{Highlighting}[]
\CommentTok{\#install.packages("ggplot2")}
\FunctionTok{library}\NormalTok{(ggplot2)}
\FunctionTok{ggplot}\NormalTok{(}\AttributeTok{data=}\NormalTok{iris, }\FunctionTok{aes}\NormalTok{(}\AttributeTok{x=}\NormalTok{Sepal.Length, }\AttributeTok{y=}\NormalTok{Sepal.Width, }\AttributeTok{color=}\NormalTok{Species)) }\SpecialCharTok{+} \FunctionTok{geom\_point}\NormalTok{(}\AttributeTok{size=}\DecValTok{4}\NormalTok{)}
\end{Highlighting}
\end{Shaded}

\pandocbounded{\includegraphics[keepaspectratio]{bookdown-demo_files/figure-latex/unnamed-chunk-32-1.pdf}}

In our activities to visualize human genomic data, we will use a library called qqman, to visualize the biological association, through a plot known as the Manhattan plot.

\subsection{Visualizing GWAS}\label{visualizing-gwas}

\begin{itemize}
\item
  The simplest definition of a GWAS is the statistical or significant association between a phenotype (trait) and a genotype. This association can also be called \emph{biological association}.
\item
  Information about association of SNPs with Huntington's Disease can be found at the Chaves 2019 Huntington's disease paper
\end{itemize}

\subsubsection{Package installation}\label{package-installation}

\begin{itemize}
\tightlist
\item
  We need to install and load package qqman
\end{itemize}

\begin{Shaded}
\begin{Highlighting}[]
\DocumentationTok{\#\# install.packages("qqman")}
\FunctionTok{library}\NormalTok{(qqman)}
\end{Highlighting}
\end{Shaded}

\begin{verbatim}
## 
\end{verbatim}

\begin{verbatim}
## For example usage please run: vignette('qqman')
\end{verbatim}

\begin{verbatim}
## 
\end{verbatim}

\begin{verbatim}
## Citation appreciated but not required:
\end{verbatim}

\begin{verbatim}
## Turner, (2018). qqman: an R package for visualizing GWAS results using Q-Q and manhattan plots. Journal of Open Source Software, 3(25), 731, https://doi.org/10.21105/joss.00731.
\end{verbatim}

\begin{verbatim}
## 
\end{verbatim}

\subsubsection{Assign text file to dataframe object}\label{assign-text-file-to-dataframe-object}

\begin{itemize}
\item
  After, we load data-frame to be visualized
\item
  Exact location of text file in your system needs to be determined
\end{itemize}

\begin{Shaded}
\begin{Highlighting}[]
\NormalTok{GWAS\_TABLE }\OtherTok{\textless{}{-}} \FunctionTok{read.table}\NormalTok{(}\StringTok{"./data/chr4.txt"}\NormalTok{, }
                         \AttributeTok{sep =} \StringTok{" "}\NormalTok{,}
                         \AttributeTok{header =}\NormalTok{ T)}
\end{Highlighting}
\end{Shaded}

\subsubsection{Get information about object with head()}\label{get-information-about-object-with-head}

\begin{itemize}
\tightlist
\item
  Use head() function to inspect the data-frame just loaded
\end{itemize}

\begin{Shaded}
\begin{Highlighting}[]
\FunctionTok{head}\NormalTok{(GWAS\_TABLE, }\AttributeTok{n =} \DecValTok{10}\NormalTok{)}
\end{Highlighting}
\end{Shaded}

\begin{verbatim}
##    CHR         SNP     BP         P
## 1    4 chr4:128096 128096 0.0133500
## 2    4 chr4:516586 516586 0.0076260
## 3    4 chr4:523979 523979 0.0024960
## 4    4 chr4:527217 527217 0.0217400
## 5    4 chr4:566177 566177 0.0008988
## 6    4 chr4:578679 578679 0.0162100
## 7    4 chr4:578790 578790 0.0103700
## 8    4 chr4:579307 579307 0.0334600
## 9    4 chr4:580259 580259 0.0190400
## 10   4 chr4:585318 585318 0.0317600
\end{verbatim}

\subsubsection{Plotting GWAS data-frame}\label{plotting-gwas-data-frame}

\begin{itemize}
\item
  In this section we use functions plot() and boxplot() to visualize data-frame
\item
  In the y access, we see genomic coordinates and in the x access, p-valoues of the biological association
\item
  Note that depending on the synthax used for plotting, R may complain
\item
  Comments not turned off
\end{itemize}

\begin{Shaded}
\begin{Highlighting}[]
\FunctionTok{plot}\NormalTok{(}\AttributeTok{data=}\NormalTok{GWAS\_TABLE, GWAS\_TABLE}\SpecialCharTok{$}\NormalTok{BP, GWAS\_TABLE}\SpecialCharTok{$}\NormalTok{P) }\DocumentationTok{\#\# R complains}
\end{Highlighting}
\end{Shaded}

\begin{verbatim}
## Warning in plot.window(...): "data" is not a graphical parameter
\end{verbatim}

\begin{verbatim}
## Warning in plot.xy(xy, type, ...): "data" is not a graphical parameter
\end{verbatim}

\begin{verbatim}
## Warning in axis(side = side, at = at, labels = labels, ...): "data" is not a
## graphical parameter
## Warning in axis(side = side, at = at, labels = labels, ...): "data" is not a
## graphical parameter
\end{verbatim}

\begin{verbatim}
## Warning in box(...): "data" is not a graphical parameter
\end{verbatim}

\begin{verbatim}
## Warning in title(...): "data" is not a graphical parameter
\end{verbatim}

\pandocbounded{\includegraphics[keepaspectratio]{bookdown-demo_files/figure-latex/3. Gwas-1.pdf}}

\begin{itemize}
\tightlist
\item
  Turn off comments
\end{itemize}

\begin{Shaded}
\begin{Highlighting}[]
\FunctionTok{plot}\NormalTok{(}\AttributeTok{data=}\NormalTok{GWAS\_TABLE, GWAS\_TABLE}\SpecialCharTok{$}\NormalTok{BP, GWAS\_TABLE}\SpecialCharTok{$}\NormalTok{P) }\DocumentationTok{\#\# R complains}
\end{Highlighting}
\end{Shaded}

\pandocbounded{\includegraphics[keepaspectratio]{bookdown-demo_files/figure-latex/unnamed-chunk-36-1.pdf}}

\subsubsection{Plotting GWAS data-frame}\label{plotting-gwas-data-frame-1}

\begin{Shaded}
\begin{Highlighting}[]
\FunctionTok{plot}\NormalTok{(GWAS\_TABLE}\SpecialCharTok{$}\NormalTok{BP, GWAS\_TABLE}\SpecialCharTok{$}\NormalTok{P) }\DocumentationTok{\#\# R does not complain}
\end{Highlighting}
\end{Shaded}

\pandocbounded{\includegraphics[keepaspectratio]{bookdown-demo_files/figure-latex/unnamed-chunk-37-1.pdf}}

\subsubsection{Boxplot}\label{boxplot-1}

\begin{itemize}
\tightlist
\item
  Boxplot
\end{itemize}

\begin{Shaded}
\begin{Highlighting}[]
\FunctionTok{boxplot}\NormalTok{(}\AttributeTok{data=}\NormalTok{GWAS\_TABLE, GWAS\_TABLE}\SpecialCharTok{$}\NormalTok{BP, GWAS\_TABLE}\SpecialCharTok{$}\NormalTok{P)}
\end{Highlighting}
\end{Shaded}

\pandocbounded{\includegraphics[keepaspectratio]{bookdown-demo_files/figure-latex/unnamed-chunk-38-1.pdf}}

\begin{itemize}
\item
  Compare the chromosomal coordinates and the values of the p-values in the boxplot above
\item
  The chunk below allows one to ommit NA values in data-frame
\end{itemize}

\begin{Shaded}
\begin{Highlighting}[]
\NormalTok{GWAS\_TABLE\_Ommit }\OtherTok{\textless{}{-}} \FunctionTok{na.omit}\NormalTok{(GWAS\_TABLE)}
\end{Highlighting}
\end{Shaded}

A seguir, criamos uma variável que armazena as posições das SNPs a serem realçadas em verde no Manhattan plot. Estas SNPs são mutações biológica e estatisticamente associadas à Doença de Huntington no gene de uma sortilina, localizada em proximidade física ao gene da proteína huntingtina mutada, a qual é a causadora medeliana (segue as leis de segregação genética de Mendel) da Doença de Huntington.

\section{Create vector to highlight genome coordinates in Manhattan plot}\label{create-vector-to-highlight-genome-coordinates-in-manhattan-plot}

\begin{Shaded}
\begin{Highlighting}[]
\NormalTok{SNP\_HIGHLIGHT }\OtherTok{\textless{}{-}} \FunctionTok{c}\NormalTok{(}\StringTok{"chr4:3043512"}\NormalTok{,}\StringTok{"chr4:3043513"}\NormalTok{,}\StringTok{"chr4:3048207"}\NormalTok{,}\StringTok{"chr4:3224216"}\NormalTok{,}
                   \StringTok{"chr4:3231772"}\NormalTok{,}\StringTok{"chr4:3233844"}\NormalTok{,}\StringTok{"chr4:3235081"}\NormalTok{,}\StringTok{"chr4:3235084"}\NormalTok{,}
                   \StringTok{"chr4:3236881"}\NormalTok{,}\StringTok{"chr4:3236883"}\NormalTok{,}\StringTok{"chr4:3241845"}\NormalTok{,}\StringTok{"chr4:3243804"}\NormalTok{,}
                   \StringTok{"chr4:3263138"}\NormalTok{,}\StringTok{"chr4:3265130"}\NormalTok{,}\StringTok{"chr4:3265710"}\NormalTok{,}\StringTok{"chr4:3314646"}\NormalTok{,}
                   \StringTok{"chr4:3380088"}\NormalTok{,}\StringTok{"chr4:3409359"}\NormalTok{,}\StringTok{"chr4:3411110"}\NormalTok{,}\StringTok{"chr4:3415336"}\NormalTok{,}
                   \StringTok{"chr4:3415378"}\NormalTok{,}\StringTok{"chr4:3438643"}\NormalTok{,}\StringTok{"chr4:3446091"}\NormalTok{,}\StringTok{"chr4:3449886"}\NormalTok{,}
                   \StringTok{"chr4:3473066"}\NormalTok{,}\StringTok{"chr4:3476809"}\NormalTok{,}\StringTok{"chr4:3480439"}\NormalTok{,}\StringTok{"chr4:3487151"}\NormalTok{,}
                   \StringTok{"chr4:3496058"}\NormalTok{,}\StringTok{"chr4:3496110"}\NormalTok{,}\StringTok{"chr4:3506933"}\NormalTok{,}\StringTok{"chr4:3508752"}\NormalTok{,}
                   \StringTok{"chr4:3510957"}\NormalTok{,}\StringTok{"chr4:3512690"}\NormalTok{,}\StringTok{"chr4:3517746"}\NormalTok{,}\StringTok{"chr4:3518190"}\NormalTok{,}
                   \StringTok{"chr4:3529671"}\NormalTok{,}\StringTok{"chr4:3532327"}\NormalTok{,}\StringTok{"chr4:3533066"}\NormalTok{,}\StringTok{"chr4:3746133"}\NormalTok{,}
                   \StringTok{"chr4:3747842"}\NormalTok{,}\StringTok{"chr4:3748134"}\NormalTok{,}\StringTok{"chr4:3765305"}\NormalTok{,}\StringTok{"chr4:3765336"}\NormalTok{,}
                   \StringTok{"chr4:3944253"}\NormalTok{,}\StringTok{"chr4:3944752"}\NormalTok{,}\StringTok{"chr4:3944888"}\NormalTok{,}\StringTok{"chr4:3946166"}\NormalTok{,}
                   \StringTok{"chr4:3946175"}\NormalTok{,}\StringTok{"chr4:3969218"}\NormalTok{,}\StringTok{"chr4:4051294"}\NormalTok{,}\StringTok{"chr4:4076788"}\NormalTok{,}
                   \StringTok{"chr4:4103104"}\NormalTok{,}\StringTok{"chr4:4103105"}\NormalTok{,}\StringTok{"chr4:4109198"}\NormalTok{,}\StringTok{"chr4:4109210"}\NormalTok{,}
                   \StringTok{"chr4:4240627"}\NormalTok{,}\StringTok{"chr4:4242705"}\NormalTok{,}\StringTok{"chr4:4243668"}\NormalTok{,}\StringTok{"chr4:4245210"}\NormalTok{,}
                   \StringTok{"chr4:4245510"}\NormalTok{,}\StringTok{"chr4:4245513"}\NormalTok{,}\StringTok{"chr4:4245591"}\NormalTok{,}\StringTok{"chr4:4245926"}\NormalTok{,}
                   \StringTok{"chr4:4245929"}\NormalTok{,}\StringTok{"chr4:4246109"}\NormalTok{,}\StringTok{"chr4:4246433"}\NormalTok{,}\StringTok{"chr4:4246453"}\NormalTok{,}
                   \StringTok{"chr4:4246457"}\NormalTok{,}\StringTok{"chr4:4246497"}\NormalTok{,}\StringTok{"chr4:4249414"}\NormalTok{,}\StringTok{"chr4:4249415"}\NormalTok{,}
                   \StringTok{"chr4:4249484"}\NormalTok{,}\StringTok{"chr4:4271623"}\NormalTok{,}\StringTok{"chr4:4275306"}\NormalTok{,}\StringTok{"chr4:4304749"}\NormalTok{,}
                   \StringTok{"chr4:4318931"}\NormalTok{,}\StringTok{"chr4:4318970"}\NormalTok{,}\StringTok{"chr4:4319564"}\NormalTok{,}\StringTok{"chr4:4319728"}\NormalTok{,}
                   \StringTok{"chr4:4319750"}\NormalTok{,}\StringTok{"chr4:4322078"}\NormalTok{,}\StringTok{"chr4:4709657"}\NormalTok{,}\StringTok{"chr4:4732282"}\NormalTok{,}
                   \StringTok{"chr4:4789635"}\NormalTok{,}\StringTok{"chr4:4822960"}\NormalTok{,}\StringTok{"chr4:4824890"}\NormalTok{,}\StringTok{"chr4:4825092"}\NormalTok{,}
                   \StringTok{"chr4:4825180"}\NormalTok{,}\StringTok{"chr4:4865316"}\NormalTok{,}\StringTok{"chr4:4865321"}\NormalTok{,}\StringTok{"chr4:5018702"}\NormalTok{,}
                   \StringTok{"chr4:5812778"}\NormalTok{,}\StringTok{"chr4:5814082"}\NormalTok{,}\StringTok{"chr4:5833660"}\NormalTok{,}\StringTok{"chr4:5833899"}\NormalTok{,}
                   \StringTok{"chr4:5835541"}\NormalTok{,}\StringTok{"chr4:5851205"}\NormalTok{,}\StringTok{"chr4:5862752"}\NormalTok{,}\StringTok{"chr4:5862938"}\NormalTok{,}
                   \StringTok{"chr4:5862943"}\NormalTok{,}\StringTok{"chr4:5901873"}\NormalTok{,}\StringTok{"chr4:5905499"}\NormalTok{,}\StringTok{"chr4:5906287"}\NormalTok{,}
                   \StringTok{"chr4:6018891"}\NormalTok{,}\StringTok{"chr4:6019046"}\NormalTok{,}\StringTok{"chr4:6020190"}\NormalTok{,}\StringTok{"chr4:6020367"}\NormalTok{,}
                   \StringTok{"chr4:6025638"}\NormalTok{,}\StringTok{"chr4:6025656"}\NormalTok{,}\StringTok{"chr4:6025766"}\NormalTok{,}\StringTok{"chr4:6026058"}\NormalTok{,}
                   \StringTok{"chr4:6083488"}\NormalTok{,}\StringTok{"chr4:6204935"}\NormalTok{,}\StringTok{"chr4:6235553"}\NormalTok{,}\StringTok{"chr4:6237142"}\NormalTok{,}
                   \StringTok{"chr4:6238466"}\NormalTok{,}\StringTok{"chr4:6239906"}\NormalTok{,}\StringTok{"chr4:6240929"}\NormalTok{,}\StringTok{"chr4:6245022"}\NormalTok{,}
                   \StringTok{"chr4:6245618"}\NormalTok{,}\StringTok{"chr4:6245732"}\NormalTok{,}\StringTok{"chr4:6245915"}\NormalTok{,}\StringTok{"chr4:6246075"}\NormalTok{,}
                   \StringTok{"chr4:6246373"}\NormalTok{,}\StringTok{"chr4:6246959"}\NormalTok{,}\StringTok{"chr4:6290594"}\NormalTok{,}\StringTok{"chr4:6292020"}\NormalTok{,}
                   \StringTok{"chr4:6294095"}\NormalTok{,}\StringTok{"chr4:6298375"}\NormalTok{,}\StringTok{"chr4:6316092"}\NormalTok{,}\StringTok{"chr4:6321396"}\NormalTok{,}
                   \StringTok{"chr4:6324647"}\NormalTok{,}\StringTok{"chr4:6324785"}\NormalTok{,}\StringTok{"chr4:6327669"}\NormalTok{,}\StringTok{"chr4:6328354"}\NormalTok{,}
                   \StringTok{"chr4:6328507"}\NormalTok{,}\StringTok{"chr4:6333130"}\NormalTok{,}\StringTok{"chr4:6333559"}\NormalTok{,}\StringTok{"chr4:6333669"}\NormalTok{,}
                   \StringTok{"chr4:6335966"}\NormalTok{,}\StringTok{"chr4:6435341"}\NormalTok{,}\StringTok{"chr4:6435486"}\NormalTok{,}\StringTok{"chr4:6435926"}\NormalTok{,}
                   \StringTok{"chr4:6437191"}\NormalTok{,}\StringTok{"chr4:6437197"}\NormalTok{,}\StringTok{"chr4:6457121"}\NormalTok{,}\StringTok{"chr4:6457131"}\NormalTok{,}
                   \StringTok{"chr4:6457132"}\NormalTok{,}\StringTok{"chr4:6568390"}\NormalTok{,}\StringTok{"chr4:6570032"}\NormalTok{,}\StringTok{"chr4:6570768"}\NormalTok{,}
                   \StringTok{"chr4:6596360"}\NormalTok{,}\StringTok{"chr4:6613252"}\NormalTok{,}\StringTok{"chr4:6613462"}\NormalTok{,}\StringTok{"chr4:6620991"}\NormalTok{,}
                   \StringTok{"chr4:6624771"}\NormalTok{,}\StringTok{"chr4:6626154"}\NormalTok{,}\StringTok{"chr4:6641969"}\NormalTok{,}\StringTok{"chr4:6642090"}\NormalTok{,}
                   \StringTok{"chr4:6644466"}\NormalTok{,}\StringTok{"chr4:6644467"}\NormalTok{,}\StringTok{"chr4:6644468"}\NormalTok{,}\StringTok{"chr4:6647889"}\NormalTok{,}
                   \StringTok{"chr4:6648300"}\NormalTok{,}\StringTok{"chr4:6662665"}\NormalTok{,}\StringTok{"chr4:6663319"}\NormalTok{,}\StringTok{"chr4:6663715"}\NormalTok{,}
                   \StringTok{"chr4:6674554"}\NormalTok{,}\StringTok{"chr4:6678553"}\NormalTok{,}\StringTok{"chr4:6678599"}\NormalTok{,}\StringTok{"chr4:6690535"}\NormalTok{,}
                   \StringTok{"chr4:6698664"}\NormalTok{,}\StringTok{"chr4:6698667"}\NormalTok{,}\StringTok{"chr4:6698706"}\NormalTok{,}\StringTok{"chr4:6720572"}\NormalTok{,}
                   \StringTok{"chr4:6911679"}\NormalTok{,}\StringTok{"chr4:6985889"}\NormalTok{,}\StringTok{"chr4:6987394"}\NormalTok{,}\StringTok{"chr4:7002344"}\NormalTok{,}
                   \StringTok{"chr4:7004495"}\NormalTok{,}\StringTok{"chr4:7004506"}\NormalTok{,}\StringTok{"chr4:7005196"}\NormalTok{,}\StringTok{"chr4:7005199"}\NormalTok{,}
                   \StringTok{"chr4:7024077"}\NormalTok{,}\StringTok{"chr4:7024398"}\NormalTok{,}\StringTok{"chr4:7029430"}\NormalTok{,}\StringTok{"chr4:7031064"}\NormalTok{,}
                   \StringTok{"chr4:7044357"}\NormalTok{,}\StringTok{"chr4:7044380"}\NormalTok{,}\StringTok{"chr4:7048842"}\NormalTok{,}\StringTok{"chr4:7052115"}\NormalTok{,}
                   \StringTok{"chr4:7055253"}\NormalTok{,}\StringTok{"chr4:7064243"}\NormalTok{,}\StringTok{"chr4:7067765"}\NormalTok{,}\StringTok{"chr4:7073187"}\NormalTok{,}
                   \StringTok{"chr4:7074027"}\NormalTok{,}\StringTok{"chr4:7677967"}\NormalTok{,}\StringTok{"chr4:7701947"}\NormalTok{,}\StringTok{"chr4:7702795"}\NormalTok{,}
                   \StringTok{"chr4:7703505"}\NormalTok{,}\StringTok{"chr4:7703807"}\NormalTok{,}\StringTok{"chr4:7704795"}\NormalTok{,}\StringTok{"chr4:7704818"}\NormalTok{,}
                   \StringTok{"chr4:7709703"}\NormalTok{,}\StringTok{"chr4:7712150"}\NormalTok{,}\StringTok{"chr4:7714490"}\NormalTok{,}\StringTok{"chr4:7733843"}\NormalTok{,}
                   \StringTok{"chr4:7735162"}\NormalTok{,}\StringTok{"chr4:7735164"}\NormalTok{,}\StringTok{"chr4:7736103"}\NormalTok{,}\StringTok{"chr4:7736112"}\NormalTok{)}
\end{Highlighting}
\end{Shaded}

\section{Biological Association Visualization: Manhattan Plot}\label{biological-association-visualization-manhattan-plot}

\begin{itemize}
\tightlist
\item
  Finally, we plot the Manhattanh plot graph
\end{itemize}

\begin{Shaded}
\begin{Highlighting}[]
\FunctionTok{manhattan}\NormalTok{(GWAS\_TABLE\_Ommit, }
          \AttributeTok{highlight =}\NormalTok{ SNP\_HIGHLIGHT, }
          \AttributeTok{annotateTop =}\NormalTok{ T, }
          \AttributeTok{annotatePval =} \FloatTok{0.20}\NormalTok{, }
          \AttributeTok{genomewideline =} \SpecialCharTok{{-}}\FunctionTok{log10}\NormalTok{(}\FloatTok{0.0001}\NormalTok{))}
\end{Highlighting}
\end{Shaded}

\pandocbounded{\includegraphics[keepaspectratio]{bookdown-demo_files/figure-latex/3. manhattan plot-1.pdf}}

\section{References}\label{references}

\section{Session Info}\label{session-info}

\begin{Shaded}
\begin{Highlighting}[]
\FunctionTok{sessionInfo}\NormalTok{()}
\end{Highlighting}
\end{Shaded}

\begin{verbatim}
## R version 4.4.1 (2024-06-14)
## Platform: aarch64-apple-darwin20
## Running under: macOS 15.7.4
## 
## Matrix products: default
## BLAS:   /Library/Frameworks/R.framework/Versions/4.4-arm64/Resources/lib/libRblas.0.dylib 
## LAPACK: /Library/Frameworks/R.framework/Versions/4.4-arm64/Resources/lib/libRlapack.dylib;  LAPACK version 3.12.0
## 
## locale:
## [1] en_US.UTF-8/en_US.UTF-8/en_US.UTF-8/C/en_US.UTF-8/en_US.UTF-8
## 
## time zone: America/Chicago
## tzcode source: internal
## 
## attached base packages:
## [1] stats     graphics  grDevices utils     datasets  methods   base     
## 
## other attached packages:
## [1] qqman_0.1.9   ggplot2_3.5.1
## 
## loaded via a namespace (and not attached):
##  [1] vctrs_0.6.5       cli_3.6.3         knitr_1.48        rlang_1.1.4      
##  [5] xfun_0.48         highr_0.11        generics_0.1.3    calibrate_1.7.7  
##  [9] labeling_0.4.3    glue_1.8.0        colorspace_2.1-1  htmltools_0.5.8.1
## [13] tinytex_0.53      scales_1.3.0      fansi_1.0.6       rmarkdown_2.28   
## [17] grid_4.4.1        evaluate_1.0.0    munsell_0.5.1     tibble_3.2.1     
## [21] MASS_7.3-61       fastmap_1.2.0     yaml_2.3.10       lifecycle_1.0.4  
## [25] bookdown_0.40     compiler_4.4.1    dplyr_1.1.4       pkgconfig_2.0.3  
## [29] rstudioapi_0.16.0 farver_2.1.2      digest_0.6.37     R6_2.5.1         
## [33] tidyselect_1.2.1  utf8_1.2.4        pillar_1.9.0      magrittr_2.0.3   
## [37] withr_3.0.1       tools_4.4.1       gtable_0.3.5
\end{verbatim}

\chapter{Processamento de dados públicos de neuroblastoma}\label{processamento-de-dados-puxfablicos-de-neuroblastoma}

\section{Origem dos dados}\label{origem-dos-dados}

Os dados analisados neste módulo têm origem na base de dados genômicos R2.

A plataforma R2 é uma ferramenta de análise e visualização genômica online de livre acesso que pode analisar uma grande coleção de dados públicos, mas também permite a análise protegida de conjuntos de dados privados, conforme indicado por Rijswijk (2024).

A R2 consiste em um banco de dados que armazena as informações genômicas, acoplado a um amplo conjunto de ferramentas para analisar/visualizar os conjuntos de dados.

Mais informações sobre a plataforma R2 pode ser encontradas no site da plataforma:

\url{https://hgserver1.amc.nl/cgi-bin/r2/main.cgi}

\hfill\break

\section{Carregando dados pre-processados de R2}\label{carregando-dados-pre-processados-de-r2}

\hfill\break

Nesta parte, vamos carregar uma dataframe contendo dados baixados de R2 em que já fizemos pre-processamento. Vamos gerar algumas visualizações para nos familiarizarmos com os dados. O objeto r2\_gse62564\_GSVA\_Metadata.rds deve ser baixado da pasta no Google Drive no seguinte link: \url{https://drive.google.com/drive/folders/1Kkos4HK2EP4ntrQOJE0mrHYlpeMa1yi4?usp=sharing}

Uma vez que o arquivo é baixado, pode-se carregá-lo no ambiente R usando a função readRDS, como feito abaixo:

\begin{Shaded}
\begin{Highlighting}[]
\NormalTok{r2\_gse62564\_GSVA\_Metadata }\OtherTok{\textless{}{-}} \FunctionTok{readRDS}\NormalTok{(}\StringTok{"./data/r2\_gse62564\_GSVA\_Metadata.rds"}\NormalTok{)}
\end{Highlighting}
\end{Shaded}

No neuroblastoma, o principal sistema de classificação de risco classifica os pacientes usando características como idade, estadiamento, histologia e status de amplificação do MYCN.

O prognóstico é a associação entre as características do tumor e a sobrevida do paciente.

O MYCN é um gene que codifica um fator de transcrição muito importante no prognóstico do tumor.

Vamos analisar a expressão do MYCN comparada pelo grupo de risco do neuroblastoma.

Primeiro, vamos ver que tipo de variável é o status do MYCN e a expressão do MYCN.

Verificar a classe da variável mycn\_status:

\begin{Shaded}
\begin{Highlighting}[]
\FunctionTok{class}\NormalTok{(r2\_gse62564\_GSVA\_Metadata}\SpecialCharTok{$}\NormalTok{mycn\_status)}
\end{Highlighting}
\end{Shaded}

\begin{verbatim}
## [1] "character"
\end{verbatim}

Verificar a classe da variável MYCN:

\begin{Shaded}
\begin{Highlighting}[]
\FunctionTok{class}\NormalTok{(r2\_gse62564\_GSVA\_Metadata}\SpecialCharTok{$}\NormalTok{MYCN)}
\end{Highlighting}
\end{Shaded}

\begin{verbatim}
## [1] "character"
\end{verbatim}

Ambos são do tipo ``character''. Para plotar, precisamos fazer com que a expressão MYCN seja do tipo numérico:

\begin{Shaded}
\begin{Highlighting}[]
\NormalTok{r2\_gse62564\_GSVA\_Metadata}\SpecialCharTok{$}\NormalTok{MYCN }\OtherTok{\textless{}{-}} \FunctionTok{as.numeric}\NormalTok{(}
\NormalTok{  r2\_gse62564\_GSVA\_Metadata}\SpecialCharTok{$}\NormalTok{MYCN)}
\end{Highlighting}
\end{Shaded}

Verificar novamente a classe da variável MYCN:

\begin{Shaded}
\begin{Highlighting}[]
\FunctionTok{class}\NormalTok{(r2\_gse62564\_GSVA\_Metadata}\SpecialCharTok{$}\NormalTok{MYCN)}
\end{Highlighting}
\end{Shaded}

\begin{verbatim}
## [1] "numeric"
\end{verbatim}

Agora a classe de MYCN é um vetor numérico.

\section{Primeiros gráficos com dados de neuroblastoma}\label{primeiros-gruxe1ficos-com-dados-de-neuroblastoma}

\subsection{Mycn\_status vs Expressão de MYCN}\label{mycn_status-vs-expressuxe3o-de-mycn}

Vamos plotar uma figura com a variável mycn\_status no eixo X e expressão de MYCN no eixo Y.

Vamos tentar ver se o ggplot, que usamos antes, considera a variável mycn\_status categórica como ela é agora.

\begin{Shaded}
\begin{Highlighting}[]
\FunctionTok{library}\NormalTok{(ggplot2)}
\FunctionTok{ggplot}\NormalTok{(r2\_gse62564\_GSVA\_Metadata, }\FunctionTok{aes}\NormalTok{(}\AttributeTok{x=}\NormalTok{mycn\_status, }\AttributeTok{y=}\NormalTok{MYCN, }
                                      \AttributeTok{fill=}\NormalTok{mycn\_status)) }\SpecialCharTok{+} 
  \FunctionTok{geom\_boxplot}\NormalTok{()}
\end{Highlighting}
\end{Shaded}

\pandocbounded{\includegraphics[keepaspectratio]{bookdown-demo_files/figure-latex/unnamed-chunk-46-1.pdf}}

Vemos que existem três categorias de mycn\_status: mycn\_amp, mycn\_nonamp e nas. mycn\_amp significa amplificação de MYCN, um fator de mau prognóstico. mycn\_nonamp significa ausência de amplificação de MYCN. NA significa que não há informações sobre o status de MYCN. Vemos que a amplificação de MYCN está altamente associada à expressão de MYCN, no eixo Y. Por enquanto, não vamos nos preocupar nem com a formatação dos rótulos X e Y nem com os valores N/A do status de amplificação MYCN.

\subsection{Alto Risco vs expressão de MYCN}\label{alto-risco-vs-expressuxe3o-de-mycn}

Podemos plotar o status de alto risco vs.~amplificação MYCN e deixar o argumento de preenchimento como mycn\_status

\begin{Shaded}
\begin{Highlighting}[]
\FunctionTok{library}\NormalTok{(ggplot2)}
\FunctionTok{ggplot}\NormalTok{(r2\_gse62564\_GSVA\_Metadata, }\FunctionTok{aes}\NormalTok{(}\AttributeTok{x=}\NormalTok{high\_risk, }\AttributeTok{y=}\NormalTok{MYCN, }
                                      \AttributeTok{fill=}\NormalTok{mycn\_status)) }\SpecialCharTok{+} 
    \FunctionTok{geom\_boxplot}\NormalTok{()}
\end{Highlighting}
\end{Shaded}

\pandocbounded{\includegraphics[keepaspectratio]{bookdown-demo_files/figure-latex/unnamed-chunk-47-1.pdf}}

Aqui vemos que a maior parte dos pacientes de alto risco (High Risk) apresentam também amplificação de MYCN e que nesses pacientes, a expressão de MYCN é alta.

Podemos ainda querer ver apenas a expressão de MYCN no grupo de alto risco, usando o parâmetro fill=high\_risk:

\begin{Shaded}
\begin{Highlighting}[]
\FunctionTok{library}\NormalTok{(ggplot2)}
\FunctionTok{ggplot}\NormalTok{(r2\_gse62564\_GSVA\_Metadata, }\FunctionTok{aes}\NormalTok{(}\AttributeTok{x=}\NormalTok{high\_risk, }\AttributeTok{y=}\NormalTok{MYCN, }
                                      \AttributeTok{fill=}\NormalTok{high\_risk)) }\SpecialCharTok{+} 
    \FunctionTok{geom\_boxplot}\NormalTok{()}
\end{Highlighting}
\end{Shaded}

\pandocbounded{\includegraphics[keepaspectratio]{bookdown-demo_files/figure-latex/unnamed-chunk-48-1.pdf}}

Onde confirmamos que a expressão de MYCN é maior nos pacientes de alto risco (High Risk).

\chapter{Análise de Sobrevivência}\label{anuxe1lise-de-sobrevivuxeancia}

Na análise de sobrevivência estudamos como varia a probabilidade de sobrevivência dos pacientes em função do tempo transcorrido desde o início de um estudo. Na atividade proposta para este módulo, vamos avaliar como é a probabilidade de sobrevivência de pacientes de neuroblastoma em função do status de alto risco (se um paciente é de alto risco ou baixo risco). Para dar prosseguimento, carregamos a dataframe de neuroblastoma no ambiente R, como fizemos no módulo anterior:

\begin{Shaded}
\begin{Highlighting}[]
\NormalTok{r2\_gse62564\_GSVA\_Metadata }\OtherTok{\textless{}{-}} \FunctionTok{readRDS}\NormalTok{( }\StringTok{"./data/r2\_gse62564\_GSVA\_Metadata.rds"}\NormalTok{)}
\end{Highlighting}
\end{Shaded}

\section{Processamento de variáveis para permitir a construção de gráficos}\label{processamento-de-variuxe1veis-para-permitir-a-construuxe7uxe3o-de-gruxe1ficos}

Precisamos definir o status do OS (Overall Survival = sobrevivência geral) como numérico, em vez de um caractere. Fazemos o mesmo para o status EFS (Event Free Survival = sobrevivência livre de evento).

Usamos o método ifelse para construir os\_bin\_numeric a partir de os\_bin, introduzindo números ao invés da palavra evento e não evento. No caso do Overall Survival (OS), a morte do paciente é o evento.

Para usar o comando a seguir, precisamos de usar, antes, uma biblioteca do R chamada maditr:

\begin{Shaded}
\begin{Highlighting}[]
\FunctionTok{library}\NormalTok{(maditr)}
\end{Highlighting}
\end{Shaded}

\begin{verbatim}
## 
## To select rows from data: rows(mtcars, am==0)
\end{verbatim}

\begin{verbatim}
## 
## Attaching package: 'maditr'
\end{verbatim}

\begin{verbatim}
## The following object is masked from 'package:base':
## 
##     sort_by
\end{verbatim}

Agora podemos usar o operador pipe (\%\textgreater\%), que toma o resultado de um comando e passa para o próximo comando:

\begin{Shaded}
\begin{Highlighting}[]
\NormalTok{r2\_gse62564\_GSVA\_Metadata }\OtherTok{\textless{}{-}}\NormalTok{ r2\_gse62564\_GSVA\_Metadata }\SpecialCharTok{\%\textgreater{}\%}
\NormalTok{  dplyr}\SpecialCharTok{::}\FunctionTok{mutate}\NormalTok{(}\AttributeTok{os\_bin\_numeric =}
                  \FunctionTok{ifelse}\NormalTok{(r2\_gse62564\_GSVA\_Metadata}\SpecialCharTok{$}\NormalTok{os\_bin }\SpecialCharTok{==} \StringTok{"event"}\NormalTok{, }\StringTok{"1"}\NormalTok{, }\StringTok{"0"}\NormalTok{))}
\end{Highlighting}
\end{Shaded}

Agora, fazemos os\_bin\_numeric numérico:

\begin{Shaded}
\begin{Highlighting}[]
\NormalTok{r2\_gse62564\_GSVA\_Metadata[}\StringTok{"os\_bin\_numeric"}\NormalTok{] }\OtherTok{\textless{}{-}} \FunctionTok{as.numeric}\NormalTok{(r2\_gse62564\_GSVA\_Metadata}\SpecialCharTok{$}\NormalTok{os\_bin\_numeric)}
\end{Highlighting}
\end{Shaded}

Fazemos o mesmo para EFS:

\begin{Shaded}
\begin{Highlighting}[]
\DocumentationTok{\#\# EFS}
\NormalTok{r2\_gse62564\_GSVA\_Metadata }\OtherTok{\textless{}{-}}\NormalTok{ r2\_gse62564\_GSVA\_Metadata }\SpecialCharTok{\%\textgreater{}\%}
\NormalTok{  dplyr}\SpecialCharTok{::}\FunctionTok{mutate}\NormalTok{(}\AttributeTok{efs\_bin\_numeric =}
                  \FunctionTok{ifelse}\NormalTok{(r2\_gse62564\_GSVA\_Metadata}\SpecialCharTok{$}\NormalTok{efs\_bin }\SpecialCharTok{==} \StringTok{"event"}\NormalTok{, }\StringTok{"1"}\NormalTok{, }\StringTok{"0"}\NormalTok{))}
\end{Highlighting}
\end{Shaded}

Agora, fazemos efs\_bin\_numeric, numérico:

\begin{Shaded}
\begin{Highlighting}[]
\NormalTok{r2\_gse62564\_GSVA\_Metadata[}\StringTok{"efs\_bin\_numeric"}\NormalTok{] }\OtherTok{\textless{}{-}} \FunctionTok{as.numeric}\NormalTok{(r2\_gse62564\_GSVA\_Metadata}\SpecialCharTok{$}\NormalTok{efs\_bin\_numeric)}
\end{Highlighting}
\end{Shaded}

\section{Fazendo a análise de sobrevivência}\label{fazendo-a-anuxe1lise-de-sobrevivuxeancia}

\subsection{Construção da curva de sobrevivência}\label{construuxe7uxe3o-da-curva-de-sobrevivuxeancia}

Usamos a função survfit do pacote em R, survival, para construir o objeto fit\_os\_high\_risk, ajustando o modelo:

\begin{Shaded}
\begin{Highlighting}[]
\FunctionTok{library}\NormalTok{(survival)}
\NormalTok{fit\_os\_high\_risk }\OtherTok{\textless{}{-}} \FunctionTok{survfit}\NormalTok{(}\FunctionTok{Surv}\NormalTok{(}\FunctionTok{as.numeric}\NormalTok{(os\_day), os\_bin\_numeric) }\SpecialCharTok{\textasciitilde{}}\NormalTok{ r2\_gse62564\_GSVA\_Metadata}\SpecialCharTok{$}\NormalTok{high\_risk,}
                   \AttributeTok{data =}\NormalTok{ r2\_gse62564\_GSVA\_Metadata)}
\end{Highlighting}
\end{Shaded}

Observe que os\_day é uma variável contendo o número de dias até o evento de interesse.

\subsection{Plotagem da curva de sobrevivência}\label{plotagem-da-curva-de-sobrevivuxeancia}

Usamos a função ggsurvplot do pacote em R survminer para plotar a curva de sobrevivência:

\begin{Shaded}
\begin{Highlighting}[]
\FunctionTok{library}\NormalTok{(survminer)}
\end{Highlighting}
\end{Shaded}

\begin{verbatim}
## Loading required package: ggpubr
\end{verbatim}

\begin{verbatim}
## 
## Attaching package: 'survminer'
\end{verbatim}

\begin{verbatim}
## The following object is masked from 'package:survival':
## 
##     myeloma
\end{verbatim}

\begin{Shaded}
\begin{Highlighting}[]
\FunctionTok{ggsurvplot}\NormalTok{(fit\_os\_high\_risk, }
           \AttributeTok{data =}\NormalTok{ r2\_gse62564\_GSVA\_Metadata, }
           \AttributeTok{risk.table =} \ConstantTok{TRUE}\NormalTok{,}
           \CommentTok{\# risk.table = "abs\_pct",}
           \AttributeTok{pval =} \ConstantTok{TRUE}\NormalTok{, }
           \AttributeTok{pval.coord =} \FunctionTok{c}\NormalTok{(}\DecValTok{2000}\NormalTok{, }\DecValTok{1}\NormalTok{),}
\NormalTok{           )}
\end{Highlighting}
\end{Shaded}

\pandocbounded{\includegraphics[keepaspectratio]{bookdown-demo_files/figure-latex/unnamed-chunk-56-1.pdf}}

\subsection{Hipóxia no microambiente tumoral}\label{hipuxf3xia-no-microambiente-tumoral}

Vamos agora procurar estabelecer uma relação entre hipóxia e neuroblastoma de alto risco. Chama-se hipóxia à condição de baixa concentração de oxigênio que se estabelece no microambiente do tumor à medida que o câncer progride.

HIF1A é um fator de transcrição sensível aos níveis de oxigênio. Precisamos tornar a coluna HIF1A (variável), numérica, como fizemos com MYCN acima:

\begin{Shaded}
\begin{Highlighting}[]
\NormalTok{r2\_gse62564\_GSVA\_Metadata}\SpecialCharTok{$}\NormalTok{HIF1A }\OtherTok{\textless{}{-}} \FunctionTok{as.numeric}\NormalTok{(r2\_gse62564\_GSVA\_Metadata}\SpecialCharTok{$}\NormalTok{HIF1A)}
\end{Highlighting}
\end{Shaded}

Compare as médias de HIF1A entre os grupos de alto risco usando ggplot2:

\begin{Shaded}
\begin{Highlighting}[]
\FunctionTok{ggplot}\NormalTok{(r2\_gse62564\_GSVA\_Metadata, }\FunctionTok{aes}\NormalTok{(}\AttributeTok{x=}\NormalTok{high\_risk, }\AttributeTok{y=}\NormalTok{HIF1A, }\AttributeTok{fill=}\NormalTok{high\_risk)) }\SpecialCharTok{+} 
    \FunctionTok{geom\_boxplot}\NormalTok{()}\SpecialCharTok{+}
  \FunctionTok{stat\_compare\_means}\NormalTok{(}\AttributeTok{method =} \StringTok{"t.test"}\NormalTok{, }\AttributeTok{label.x =} \FloatTok{1.4}\NormalTok{)}\SpecialCharTok{+}
  \FunctionTok{theme\_linedraw}\NormalTok{()}
\end{Highlighting}
\end{Shaded}

\pandocbounded{\includegraphics[keepaspectratio]{bookdown-demo_files/figure-latex/unnamed-chunk-58-1.pdf}}

Vemos que para este conjunto de dados, a expressão de HIF1A é maior no grupo de alto risco, o que está de acordo com a noção de que aumento da hipóxia leva a maior malignidade na progressão tumoral.

\subsection{Análise de sobrevivência do gene da hipóxia}\label{anuxe1lise-de-sobrevivuxeancia-do-gene-da-hipuxf3xia}

Como vemos que indivíduos de alto risco apresentam maior expressão de HIF1A, vamos analisar também a análise de sobrevivência de HIF1A. Para isso vamos separar em dois grupos de pacientes, baseados na expressão do HIF1A. Construindo a variável HIF1A em dois quantis, um com alta expressão de HIF1A e um com baixa expressão de HIF1A. Para isso vamos usar a função ntile() do pacote tidyverse:

\begin{Shaded}
\begin{Highlighting}[]
\FunctionTok{library}\NormalTok{(tidyverse)}
\end{Highlighting}
\end{Shaded}

\begin{verbatim}
## -- Attaching core tidyverse packages ------------------------ tidyverse 2.0.0 --
## v dplyr     1.1.4     v readr     2.1.5
## v forcats   1.0.0     v stringr   1.5.1
## v lubridate 1.9.3     v tibble    3.2.1
## v purrr     1.0.2     v tidyr     1.3.1
## -- Conflicts ------------------------------------------ tidyverse_conflicts() --
## x dplyr::between()   masks maditr::between()
## x dplyr::coalesce()  masks maditr::coalesce()
## x readr::cols()      masks maditr::cols()
## x dplyr::filter()    masks stats::filter()
## x dplyr::first()     masks maditr::first()
## x dplyr::lag()       masks stats::lag()
## x dplyr::last()      masks maditr::last()
## x purrr::transpose() masks maditr::transpose()
## i Use the conflicted package (<http://conflicted.r-lib.org/>) to force all conflicts to become errors
\end{verbatim}

\begin{Shaded}
\begin{Highlighting}[]
\NormalTok{r2\_gse62564\_GSVA\_Metadata }\OtherTok{\textless{}{-}}\NormalTok{ r2\_gse62564\_GSVA\_Metadata }\SpecialCharTok{\%\textgreater{}\%}
    \FunctionTok{mutate}\NormalTok{(}\AttributeTok{quantile =} \FunctionTok{ntile}\NormalTok{(HIF1A, }\DecValTok{2}\NormalTok{))}
\end{Highlighting}
\end{Shaded}

Fazendo da variável quantil HIF1A um fator

\begin{Shaded}
\begin{Highlighting}[]
\NormalTok{r2\_gse62564\_GSVA\_Metadata}\SpecialCharTok{$}\NormalTok{quantile }\OtherTok{\textless{}{-}} \FunctionTok{as.factor}\NormalTok{(r2\_gse62564\_GSVA\_Metadata}\SpecialCharTok{$}\NormalTok{quantile )}
\end{Highlighting}
\end{Shaded}

Gráfico de dispersão HALLMARK HYPOXIA vs.~HIF1A

\begin{Shaded}
\begin{Highlighting}[]
\NormalTok{r2\_gse62564\_GSVA\_Metadata}\SpecialCharTok{$}\NormalTok{HALLMARK\_HYPOXIA }\OtherTok{\textless{}{-}} \FunctionTok{as.numeric}\NormalTok{(r2\_gse62564\_GSVA\_Metadata}\SpecialCharTok{$}\NormalTok{HALLMARK\_HYPOXIA)}

\FunctionTok{qplot}\NormalTok{(HALLMARK\_HYPOXIA, HIF1A, }
      \AttributeTok{data =}\NormalTok{ r2\_gse62564\_GSVA\_Metadata,}
      \AttributeTok{colour=}\NormalTok{quantile,}
      \AttributeTok{ylab =} \StringTok{"HIF1A"}\NormalTok{,}
      \AttributeTok{xlab =} \StringTok{"Hallmark Hypoxia"}\NormalTok{)}
\end{Highlighting}
\end{Shaded}

\begin{verbatim}
## Warning: `qplot()` was deprecated in ggplot2 3.4.0.
## This warning is displayed once every 8 hours.
## Call `lifecycle::last_lifecycle_warnings()` to see where this warning was
## generated.
\end{verbatim}

\pandocbounded{\includegraphics[keepaspectratio]{bookdown-demo_files/figure-latex/unnamed-chunk-61-1.pdf}}

\begin{Shaded}
\begin{Highlighting}[]
\FunctionTok{library}\NormalTok{(ggplot2)}
\CommentTok{\# Basic scatter plot}
\FunctionTok{ggplot}\NormalTok{(r2\_gse62564\_GSVA\_Metadata, }\FunctionTok{aes}\NormalTok{(}\AttributeTok{x=}\NormalTok{HALLMARK\_HYPOXIA, }\AttributeTok{y=}\NormalTok{HIF1A, }\AttributeTok{color=}\NormalTok{quantile)) }\SpecialCharTok{+} \FunctionTok{geom\_point}\NormalTok{()}
\end{Highlighting}
\end{Shaded}

\pandocbounded{\includegraphics[keepaspectratio]{bookdown-demo_files/figure-latex/unnamed-chunk-61-2.pdf}}

Vemos que o quantil 2 tem valores de expressão de HIF1A maiores que os valores do quantil 1.

Construímos o objeto de sobrevivência ajustado com a variável quantil HIF1A:

\begin{Shaded}
\begin{Highlighting}[]
\FunctionTok{library}\NormalTok{(survival)}
\NormalTok{fit\_os\_HIF1A }\OtherTok{\textless{}{-}} \FunctionTok{survfit}\NormalTok{(}\FunctionTok{Surv}\NormalTok{(}\FunctionTok{as.numeric}\NormalTok{(os\_day), os\_bin\_numeric) }\SpecialCharTok{\textasciitilde{}}\NormalTok{ r2\_gse62564\_GSVA\_Metadata}\SpecialCharTok{$}\NormalTok{quantile,}
                   \AttributeTok{data =}\NormalTok{ r2\_gse62564\_GSVA\_Metadata)}
\end{Highlighting}
\end{Shaded}

E finalmente traçamos a curva de sobrevivência do gráfico do Overall Survival:

\begin{Shaded}
\begin{Highlighting}[]
\FunctionTok{library}\NormalTok{(survminer)}
\FunctionTok{ggsurvplot}\NormalTok{(fit\_os\_HIF1A, }
           \AttributeTok{data =}\NormalTok{ r2\_gse62564\_GSVA\_Metadata, }
           \AttributeTok{risk.table =} \ConstantTok{TRUE}\NormalTok{,}
           \CommentTok{\# risk.table = "abs\_pct",}
           \AttributeTok{pval =} \ConstantTok{TRUE}\NormalTok{, }
           \AttributeTok{pval.coord =} \FunctionTok{c}\NormalTok{(}\DecValTok{4000}\NormalTok{, }\DecValTok{1}\NormalTok{),}
\NormalTok{           )}
\end{Highlighting}
\end{Shaded}

\pandocbounded{\includegraphics[keepaspectratio]{bookdown-demo_files/figure-latex/unnamed-chunk-63-1.pdf}}

Observamos que os maiores valores na expressão de HIF1A (segundo quantil ou quantil 2. No inglês é second quantile) levam a menor probabilidade de sobrevivência. Vemos que no dataset, o segundo quantil tem menores probabilidades de sobrevivência.

\subsubsection{Event-free survival}\label{event-free-survival}

Construímos a variável quantil HIF1A:

\begin{Shaded}
\begin{Highlighting}[]
\FunctionTok{library}\NormalTok{(tidyverse)}
\NormalTok{r2\_gse62564\_GSVA\_Metadata }\OtherTok{\textless{}{-}}\NormalTok{ r2\_gse62564\_GSVA\_Metadata }\SpecialCharTok{\%\textgreater{}\%}
    \FunctionTok{mutate}\NormalTok{(}\AttributeTok{quantile =} \FunctionTok{ntile}\NormalTok{(HIF1A, }\DecValTok{2}\NormalTok{))}
\end{Highlighting}
\end{Shaded}

Construímos objeto de sobrevivência EFS, \textbf{fit\_efs\_HIF1A}

\begin{Shaded}
\begin{Highlighting}[]
\FunctionTok{library}\NormalTok{(survival)}
\NormalTok{fit\_efs\_HIF1A }\OtherTok{\textless{}{-}} \FunctionTok{survfit}\NormalTok{(}\FunctionTok{Surv}\NormalTok{(}\FunctionTok{as.numeric}\NormalTok{(efs\_day), efs\_bin\_numeric) }\SpecialCharTok{\textasciitilde{}}\NormalTok{ r2\_gse62564\_GSVA\_Metadata}\SpecialCharTok{$}\NormalTok{quantile,}
                   \AttributeTok{data =}\NormalTok{ r2\_gse62564\_GSVA\_Metadata)}
\end{Highlighting}
\end{Shaded}

Curva de sobrevivência do gráfico EFS, usando a função \textbf{ggsurvplot}

\begin{Shaded}
\begin{Highlighting}[]
\FunctionTok{library}\NormalTok{(survminer)}
\FunctionTok{ggsurvplot}\NormalTok{(fit\_efs\_HIF1A, }
           \AttributeTok{data =}\NormalTok{ r2\_gse62564\_GSVA\_Metadata, }
           \AttributeTok{risk.table =} \ConstantTok{TRUE}\NormalTok{,}
           \CommentTok{\# risk.table = "abs\_pct",}
           \AttributeTok{pval =} \ConstantTok{TRUE}\NormalTok{, }
           \AttributeTok{pval.coord =} \FunctionTok{c}\NormalTok{(}\DecValTok{4000}\NormalTok{, }\DecValTok{1}\NormalTok{),}
\NormalTok{           )}
\end{Highlighting}
\end{Shaded}

\pandocbounded{\includegraphics[keepaspectratio]{bookdown-demo_files/figure-latex/unnamed-chunk-66-1.pdf}}

\subsection{Relação do HIF1A com a idade no diagnóstico}\label{relauxe7uxe3o-do-hif1a-com-a-idade-no-diagnuxf3stico}

Precisamos criar o objeto r2\_gse62564\_GSVA\_Metadata\_Spaces\_survival

A função compare\_means não consegue lidar com espaços no nome da variável de exaustão de células T.

Teremos que usar um truque para lidar com os espaços em ``T cell exhaustion''.

Usaremos o pacote stringr para inserir pontos (.) substituindo os espaços no dataframe r2\_gse62564\_GSVA\_Metadata.

\begin{Shaded}
\begin{Highlighting}[]
\NormalTok{r2\_gse62564\_GSVA\_Metadata\_Spaces }\OtherTok{\textless{}{-}}\NormalTok{ r2\_gse62564\_GSVA\_Metadata}
\FunctionTok{library}\NormalTok{(stringr)}
\FunctionTok{names}\NormalTok{(r2\_gse62564\_GSVA\_Metadata\_Spaces)}\OtherTok{\textless{}{-}}\FunctionTok{str\_replace\_all}\NormalTok{(}\FunctionTok{names}\NormalTok{(r2\_gse62564\_GSVA\_Metadata\_Spaces), }\FunctionTok{c}\NormalTok{(}\StringTok{" "} \OtherTok{=} \StringTok{"."}\NormalTok{ , }\StringTok{","} \OtherTok{=} \StringTok{""}\NormalTok{ ))}
\end{Highlighting}
\end{Shaded}

Criando o objeto r2\_gse62564\_GSVA\_Metadata\_Spaces\_survival:

\begin{Shaded}
\begin{Highlighting}[]
\NormalTok{r2\_gse62564\_GSVA\_Metadata\_Spaces\_survival }\OtherTok{\textless{}{-}}\NormalTok{ r2\_gse62564\_GSVA\_Metadata\_Spaces}
\end{Highlighting}
\end{Shaded}

Fazendo age\_at\_diagnosis numérico para plotar os quantis

\begin{Shaded}
\begin{Highlighting}[]
\NormalTok{r2\_gse62564\_GSVA\_Metadata\_Spaces\_survival}\SpecialCharTok{$}\NormalTok{age\_at\_diagnosis }\OtherTok{\textless{}{-}} \FunctionTok{as.numeric}\NormalTok{(r2\_gse62564\_GSVA\_Metadata\_Spaces\_survival}\SpecialCharTok{$}\NormalTok{age\_at\_diagnosis)}
\end{Highlighting}
\end{Shaded}

Verificando a idade máxima no diagnóstico

\begin{Shaded}
\begin{Highlighting}[]
\FunctionTok{max}\NormalTok{(r2\_gse62564\_GSVA\_Metadata\_Spaces\_survival}\SpecialCharTok{$}\NormalTok{age\_at\_diagnosis)}
\end{Highlighting}
\end{Shaded}

\begin{verbatim}
## [1] 8983
\end{verbatim}

\begin{Shaded}
\begin{Highlighting}[]
\FunctionTok{min}\NormalTok{(r2\_gse62564\_GSVA\_Metadata\_Spaces\_survival}\SpecialCharTok{$}\NormalTok{age\_at\_diagnosis)}
\end{Highlighting}
\end{Shaded}

\begin{verbatim}
## [1] 0
\end{verbatim}

Construir quantis para a idade no diagnóstico (age\_qtls)

\begin{Shaded}
\begin{Highlighting}[]
\NormalTok{r2\_gse62564\_GSVA\_Metadata\_Spaces\_survival\_age }\OtherTok{\textless{}{-}}\NormalTok{ r2\_gse62564\_GSVA\_Metadata\_Spaces\_survival }\SpecialCharTok{\%\textgreater{}\%}
    \FunctionTok{mutate}\NormalTok{(}\AttributeTok{age\_qtls =} \FunctionTok{ntile}\NormalTok{(age\_at\_diagnosis, }\DecValTok{3}\NormalTok{))}
\end{Highlighting}
\end{Shaded}

Faça age\_qtls como fator

\begin{Shaded}
\begin{Highlighting}[]
\NormalTok{r2\_gse62564\_GSVA\_Metadata\_Spaces\_survival\_age}\SpecialCharTok{$}\NormalTok{age\_qtls }\OtherTok{\textless{}{-}} \FunctionTok{as.factor}\NormalTok{(r2\_gse62564\_GSVA\_Metadata\_Spaces\_survival\_age}\SpecialCharTok{$}\NormalTok{age\_qtls)}
\end{Highlighting}
\end{Shaded}

Gráfico age\_qtls vs.~HIF1A

\begin{Shaded}
\begin{Highlighting}[]
\FunctionTok{ggplot}\NormalTok{(r2\_gse62564\_GSVA\_Metadata\_Spaces\_survival\_age, }\FunctionTok{aes}\NormalTok{(}\AttributeTok{x=}\NormalTok{age\_qtls, }\AttributeTok{y=}\NormalTok{HIF1A, }\AttributeTok{fill=}\NormalTok{high\_risk)) }\SpecialCharTok{+} 
    \FunctionTok{geom\_boxplot}\NormalTok{()}\SpecialCharTok{+}
  \FunctionTok{stat\_compare\_means}\NormalTok{(}\AttributeTok{method =} \StringTok{"t.test"}\NormalTok{)}\SpecialCharTok{+}
  \FunctionTok{theme\_linedraw}\NormalTok{()}
\end{Highlighting}
\end{Shaded}

\pandocbounded{\includegraphics[keepaspectratio]{bookdown-demo_files/figure-latex/unnamed-chunk-74-1.pdf}}

\chapter{Accessing the Raça/Cor Variable in DATASUS}\label{accessing-the-rauxe7acor-variable-in-datasus}

\section{Load libraries}\label{load-libraries}

\hfill\break

\begin{Shaded}
\begin{Highlighting}[]
\FunctionTok{library}\NormalTok{(microdatasus)}
\FunctionTok{library}\NormalTok{(dplyr)}
\FunctionTok{library}\NormalTok{(tidyverse)}
\FunctionTok{library}\NormalTok{(maditr)}
\FunctionTok{library}\NormalTok{(kableExtra)}
\end{Highlighting}
\end{Shaded}

\hfill\break

\section{Constrir dataframe (df) com identificação racial}\label{constrir-dataframe-df-com-identificauxe7uxe3o-racial}

\hfill\break

\begin{Shaded}
\begin{Highlighting}[]
\NormalTok{racacor\_df }\OtherTok{\textless{}{-}} \FunctionTok{readRDS}\NormalTok{(}\StringTok{"./data/racacor\_df.rds"}\NormalTok{)}
\NormalTok{racacor\_df}
\end{Highlighting}
\end{Shaded}

\begin{verbatim}
##   RACACOR       n raca_cor_factor
## 1       1 6794687          branca
## 2       4 5295315           parda
## 3       2 1127508           preta
## 4    <NA>  364106            <NA>
## 5       3   74842         amarela
## 6       5   46782        indígena
## 7       9       5        ignorado
\end{verbatim}

\hfill\break

\section{Transformar variável Raça/Cor}\label{transformar-variuxe1vel-rauxe7acor}

Transforma coluna com caracteres em nomes de raça/cor

\begin{Shaded}
\begin{Highlighting}[]
\NormalTok{racacor\_df }\OtherTok{\textless{}{-}}\NormalTok{ racacor\_df }\SpecialCharTok{\%\textgreater{}\%}
\NormalTok{  dplyr}\SpecialCharTok{::}\FunctionTok{mutate}\NormalTok{(}\AttributeTok{raca\_cor\_factor =} \FunctionTok{ifelse}\NormalTok{(racacor\_df}\SpecialCharTok{$}\NormalTok{RACACOR }\SpecialCharTok{==} \StringTok{"1"}\NormalTok{, }\StringTok{"branca"}\NormalTok{, }
                                         \FunctionTok{ifelse}\NormalTok{(racacor\_df}\SpecialCharTok{$}\NormalTok{RACACOR }\SpecialCharTok{==} \StringTok{"2"}\NormalTok{, }\StringTok{"preta"}\NormalTok{,}
                                                \FunctionTok{ifelse}\NormalTok{(racacor\_df}\SpecialCharTok{$}\NormalTok{RACACOR }\SpecialCharTok{==} \StringTok{"3"}\NormalTok{, }\StringTok{"amarela"}\NormalTok{,}
                                                       \FunctionTok{ifelse}\NormalTok{(racacor\_df}\SpecialCharTok{$}\NormalTok{RACACOR }\SpecialCharTok{==} \StringTok{"4"}\NormalTok{, }\StringTok{"parda"}\NormalTok{, }
                                                              \FunctionTok{ifelse}\NormalTok{(racacor\_df}\SpecialCharTok{$}\NormalTok{RACACOR }\SpecialCharTok{==} \StringTok{"5"}\NormalTok{, }\StringTok{"indígena"}\NormalTok{,}
                                                                     \FunctionTok{ifelse}\NormalTok{(racacor\_df}\SpecialCharTok{$}\NormalTok{RACACOR }\SpecialCharTok{==} \StringTok{"9"}\NormalTok{, }\StringTok{"ignorado"}\NormalTok{,}\StringTok{"0"}\NormalTok{)))))))}
\end{Highlighting}
\end{Shaded}

\hfill\break

Verificar objeto racacor\_df

\begin{Shaded}
\begin{Highlighting}[]
\NormalTok{racacor\_df}
\end{Highlighting}
\end{Shaded}

\begin{verbatim}
##   RACACOR       n raca_cor_factor
## 1       1 6794687          branca
## 2       4 5295315           parda
## 3       2 1127508           preta
## 4    <NA>  364106            <NA>
## 5       3   74842         amarela
## 6       5   46782        indígena
## 7       9       5        ignorado
\end{verbatim}

\hfill\break

\section{Gráfico de barras por categoria da Variável Raça/Cor}\label{gruxe1fico-de-barras-por-categoria-da-variuxe1vel-rauxe7acor}

\hfill\break

Nesta parte, poderemos estimar o número de indivíduos sob risco, considerando as categorias raciais brasileiras.

\begin{Shaded}
\begin{Highlighting}[]
\FunctionTok{library}\NormalTok{(ggplot2)}
\CommentTok{\# Basic barplot}
\FunctionTok{ggplot}\NormalTok{(}\AttributeTok{data=}\NormalTok{racacor\_df, }\FunctionTok{aes}\NormalTok{(}\AttributeTok{x=}\NormalTok{raca\_cor\_factor, }\AttributeTok{y=}\NormalTok{n)) }\SpecialCharTok{+}
  \FunctionTok{geom\_bar}\NormalTok{(}\AttributeTok{stat=}\StringTok{"identity"}\NormalTok{)}
\end{Highlighting}
\end{Shaded}

\pandocbounded{\includegraphics[keepaspectratio]{bookdown-demo_files/figure-latex/barplot-1.pdf}}

\hfill\break

\begin{Shaded}
\begin{Highlighting}[]
\NormalTok{dados\_agregados\_year }\OtherTok{\textless{}{-}} \FunctionTok{readRDS}\NormalTok{(}\StringTok{"./data/dados\_agregados\_year.rds"}\NormalTok{)}
\NormalTok{dados\_agregados\_year}
\end{Highlighting}
\end{Shaded}

\begin{verbatim}
## # A tibble: 10 x 2
##    year        n
##    <chr>   <int>
##  1 2015  1264175
##  2 2016  1309774
##  3 2017  1312663
##  4 2018  1316719
##  5 2019  1349801
##  6 2020  1207189
##  7 2021  1401033
##  8 2022  1544266
##  9 2023  1465610
## 10 2024  1532015
\end{verbatim}

\hfill\break

Remover NAs dos dados agregados por ano

\begin{Shaded}
\begin{Highlighting}[]
\NormalTok{dados\_agregados\_year }\OtherTok{\textless{}{-}}\NormalTok{ dados\_agregados\_year[ }\FunctionTok{which}\NormalTok{(}
\NormalTok{  dados\_agregados\_year}\SpecialCharTok{$}\NormalTok{year }\SpecialCharTok{!=} \StringTok{"NA"}\NormalTok{), ]}
\end{Highlighting}
\end{Shaded}

\section{Mortalidade Anual Total}\label{mortalidade-anual-total}

\hfill\break

\begin{Shaded}
\begin{Highlighting}[]
\FunctionTok{library}\NormalTok{(ggplot2)}
\CommentTok{\# Basic line plot with points}
\FunctionTok{ggplot}\NormalTok{(}\AttributeTok{data=}\NormalTok{dados\_agregados\_year, }\FunctionTok{aes}\NormalTok{(}\AttributeTok{x=}\NormalTok{year, }\AttributeTok{y=}\NormalTok{n, }\AttributeTok{group=}\DecValTok{1}\NormalTok{)) }\SpecialCharTok{+}
  \FunctionTok{geom\_line}\NormalTok{()}\SpecialCharTok{+}
  \FunctionTok{geom\_point}\NormalTok{()}
\end{Highlighting}
\end{Shaded}

\pandocbounded{\includegraphics[keepaspectratio]{bookdown-demo_files/figure-latex/unnamed-chunk-78-1.pdf}}

\hfill\break

\subsection{Mortalidade Anual Total Usando Variável Raça/Cor}\label{mortalidade-anual-total-usando-variuxe1vel-rauxe7acor}

\hfill\break

Mortalidade anual total (figura acima) e mortalidade anual usando a variável Raça/Cor (figura abaixo) foram feitas usando \citet{STHTDA202408} como referência.

\hfill\break

Carregar dados\_appended\_melted

\begin{Shaded}
\begin{Highlighting}[]
\NormalTok{dados\_appended\_melted }\OtherTok{\textless{}{-}} \FunctionTok{readRDS}\NormalTok{(}\StringTok{"./data/dados\_appended\_melted.rds"}\NormalTok{)}
\NormalTok{dados\_appended\_melted}
\end{Highlighting}
\end{Shaded}

\begin{verbatim}
##    year variable  value
## 1  2015        1 643928
## 2  2016        1 665737
## 3  2017        1 663147
## 4  2018        1 668707
## 5  2019        1 686474
## 6  2020        1 523848
## 7  2021        1 644210
## 8  2022        1 787363
## 9  2023        1 737590
## 10 2024        1 773683
## 11 2015        2  93830
## 12 2016        2  97747
## 13 2017        2  99745
## 14 2018        2 102889
## 15 2019        2 107306
## 16 2020        2 107171
## 17 2021        2 122852
## 18 2022        2 131005
## 19 2023        2 129047
## 20 2024        2 135916
## 21 2015        3   7051
## 22 2016        3   7276
## 23 2017        3   7375
## 24 2018        3   7524
## 25 2019        3   7359
## 26 2020        3   4716
## 27 2021        3   5313
## 28 2022        3   9261
## 29 2023        3   9053
## 30 2024        3   9914
## 31 2015        4 458807
## 32 2016        4 483026
## 33 2017        4 497839
## 34 2018        4 496106
## 35 2019        4 509936
## 36 2020        4 531808
## 37 2021        4 587323
## 38 2022        4 582469
## 39 2023        4 561438
## 40 2024        4 586563
## 41 2015        5   3640
## 42 2016        5   3867
## 43 2017        5   3911
## 44 2018        5   4261
## 45 2019        5   4273
## 46 2020        5   5173
## 47 2021        5   5361
## 48 2022        5   5343
## 49 2023        5   5398
## 50 2024        5   5555
## 51 2015        9      0
## 52 2016        9      5
## 53 2017        9      0
## 54 2018        9      0
## 55 2019        9      0
## 56 2020        9      0
## 57 2021        9      0
## 58 2022        9      0
## 59 2023        9      0
## 60 2024        9      0
## 61 2015       NA  56919
## 62 2016       NA  52116
## 63 2017       NA  40646
## 64 2018       NA  37232
## 65 2019       NA  34453
## 66 2020       NA  34473
## 67 2021       NA  35974
## 68 2022       NA  28825
## 69 2023       NA  23084
## 70 2024       NA  20384
\end{verbatim}

\hfill\break

Obter variável raça/cor como palavras descritivas ao invés de números.

\begin{Shaded}
\begin{Highlighting}[]
\NormalTok{dados\_appended\_melted }\OtherTok{\textless{}{-}}\NormalTok{ dados\_appended\_melted }\SpecialCharTok{\%\textgreater{}\%}
\NormalTok{  dplyr}\SpecialCharTok{::}\FunctionTok{mutate}\NormalTok{(}\AttributeTok{raca\_cor\_factor =} \FunctionTok{ifelse}\NormalTok{(dados\_appended\_melted}\SpecialCharTok{$}\NormalTok{variable }\SpecialCharTok{==} \StringTok{"1"}\NormalTok{, }\StringTok{"branca"}\NormalTok{, }
                                         \FunctionTok{ifelse}\NormalTok{(dados\_appended\_melted}\SpecialCharTok{$}\NormalTok{variable }\SpecialCharTok{==} \StringTok{"2"}\NormalTok{, }\StringTok{"preta"}\NormalTok{,}
                                                \FunctionTok{ifelse}\NormalTok{(dados\_appended\_melted}\SpecialCharTok{$}\NormalTok{variable }\SpecialCharTok{==} \StringTok{"3"}\NormalTok{, }\StringTok{"amarela"}\NormalTok{,}
                                                       \FunctionTok{ifelse}\NormalTok{(dados\_appended\_melted}\SpecialCharTok{$}\NormalTok{variable }\SpecialCharTok{==} \StringTok{"4"}\NormalTok{, }\StringTok{"parda"}\NormalTok{, }
                                                              \FunctionTok{ifelse}\NormalTok{(dados\_appended\_melted}\SpecialCharTok{$}\NormalTok{variable }\SpecialCharTok{==} \StringTok{"5"}\NormalTok{, }\StringTok{"indígena"}\NormalTok{,}
                                                                     \FunctionTok{ifelse}\NormalTok{(dados\_appended\_melted}\SpecialCharTok{$}\NormalTok{variable }\SpecialCharTok{==} \StringTok{"9"}\NormalTok{, }\StringTok{"ignorado"}\NormalTok{,}\StringTok{"0"}\NormalTok{)))))))}
\CommentTok{\# Checar dados\_appended\_melted}
\NormalTok{dados\_appended\_melted}
\end{Highlighting}
\end{Shaded}

\begin{verbatim}
##    year variable  value raca_cor_factor
## 1  2015        1 643928          branca
## 2  2016        1 665737          branca
## 3  2017        1 663147          branca
## 4  2018        1 668707          branca
## 5  2019        1 686474          branca
## 6  2020        1 523848          branca
## 7  2021        1 644210          branca
## 8  2022        1 787363          branca
## 9  2023        1 737590          branca
## 10 2024        1 773683          branca
## 11 2015        2  93830           preta
## 12 2016        2  97747           preta
## 13 2017        2  99745           preta
## 14 2018        2 102889           preta
## 15 2019        2 107306           preta
## 16 2020        2 107171           preta
## 17 2021        2 122852           preta
## 18 2022        2 131005           preta
## 19 2023        2 129047           preta
## 20 2024        2 135916           preta
## 21 2015        3   7051         amarela
## 22 2016        3   7276         amarela
## 23 2017        3   7375         amarela
## 24 2018        3   7524         amarela
## 25 2019        3   7359         amarela
## 26 2020        3   4716         amarela
## 27 2021        3   5313         amarela
## 28 2022        3   9261         amarela
## 29 2023        3   9053         amarela
## 30 2024        3   9914         amarela
## 31 2015        4 458807           parda
## 32 2016        4 483026           parda
## 33 2017        4 497839           parda
## 34 2018        4 496106           parda
## 35 2019        4 509936           parda
## 36 2020        4 531808           parda
## 37 2021        4 587323           parda
## 38 2022        4 582469           parda
## 39 2023        4 561438           parda
## 40 2024        4 586563           parda
## 41 2015        5   3640        indígena
## 42 2016        5   3867        indígena
## 43 2017        5   3911        indígena
## 44 2018        5   4261        indígena
## 45 2019        5   4273        indígena
## 46 2020        5   5173        indígena
## 47 2021        5   5361        indígena
## 48 2022        5   5343        indígena
## 49 2023        5   5398        indígena
## 50 2024        5   5555        indígena
## 51 2015        9      0        ignorado
## 52 2016        9      5        ignorado
## 53 2017        9      0        ignorado
## 54 2018        9      0        ignorado
## 55 2019        9      0        ignorado
## 56 2020        9      0        ignorado
## 57 2021        9      0        ignorado
## 58 2022        9      0        ignorado
## 59 2023        9      0        ignorado
## 60 2024        9      0        ignorado
## 61 2015       NA  56919               0
## 62 2016       NA  52116               0
## 63 2017       NA  40646               0
## 64 2018       NA  37232               0
## 65 2019       NA  34453               0
## 66 2020       NA  34473               0
## 67 2021       NA  35974               0
## 68 2022       NA  28825               0
## 69 2023       NA  23084               0
## 70 2024       NA  20384               0
\end{verbatim}

\hfill\break

Remove NA in year

\begin{Shaded}
\begin{Highlighting}[]
\NormalTok{dados\_appended\_melted }\OtherTok{\textless{}{-}}\NormalTok{ dados\_appended\_melted[ }\FunctionTok{which}\NormalTok{(}
\NormalTok{  dados\_appended\_melted}\SpecialCharTok{$}\NormalTok{year }\SpecialCharTok{!=} \StringTok{"NA"}\NormalTok{), ]}
\end{Highlighting}
\end{Shaded}

\hfill\break

Remove 0 in raça/cor

\begin{Shaded}
\begin{Highlighting}[]
\NormalTok{dados\_appended\_melted }\OtherTok{\textless{}{-}}\NormalTok{ dados\_appended\_melted[ }\FunctionTok{which}\NormalTok{(}
\NormalTok{  dados\_appended\_melted}\SpecialCharTok{$}\NormalTok{raca\_cor\_factor }\SpecialCharTok{!=} \StringTok{"0"}\NormalTok{), ]}

\NormalTok{dados\_appended\_melted}
\end{Highlighting}
\end{Shaded}

\begin{verbatim}
##    year variable  value raca_cor_factor
## 1  2015        1 643928          branca
## 2  2016        1 665737          branca
## 3  2017        1 663147          branca
## 4  2018        1 668707          branca
## 5  2019        1 686474          branca
## 6  2020        1 523848          branca
## 7  2021        1 644210          branca
## 8  2022        1 787363          branca
## 9  2023        1 737590          branca
## 10 2024        1 773683          branca
## 11 2015        2  93830           preta
## 12 2016        2  97747           preta
## 13 2017        2  99745           preta
## 14 2018        2 102889           preta
## 15 2019        2 107306           preta
## 16 2020        2 107171           preta
## 17 2021        2 122852           preta
## 18 2022        2 131005           preta
## 19 2023        2 129047           preta
## 20 2024        2 135916           preta
## 21 2015        3   7051         amarela
## 22 2016        3   7276         amarela
## 23 2017        3   7375         amarela
## 24 2018        3   7524         amarela
## 25 2019        3   7359         amarela
## 26 2020        3   4716         amarela
## 27 2021        3   5313         amarela
## 28 2022        3   9261         amarela
## 29 2023        3   9053         amarela
## 30 2024        3   9914         amarela
## 31 2015        4 458807           parda
## 32 2016        4 483026           parda
## 33 2017        4 497839           parda
## 34 2018        4 496106           parda
## 35 2019        4 509936           parda
## 36 2020        4 531808           parda
## 37 2021        4 587323           parda
## 38 2022        4 582469           parda
## 39 2023        4 561438           parda
## 40 2024        4 586563           parda
## 41 2015        5   3640        indígena
## 42 2016        5   3867        indígena
## 43 2017        5   3911        indígena
## 44 2018        5   4261        indígena
## 45 2019        5   4273        indígena
## 46 2020        5   5173        indígena
## 47 2021        5   5361        indígena
## 48 2022        5   5343        indígena
## 49 2023        5   5398        indígena
## 50 2024        5   5555        indígena
## 51 2015        9      0        ignorado
## 52 2016        9      5        ignorado
## 53 2017        9      0        ignorado
## 54 2018        9      0        ignorado
## 55 2019        9      0        ignorado
## 56 2020        9      0        ignorado
## 57 2021        9      0        ignorado
## 58 2022        9      0        ignorado
## 59 2023        9      0        ignorado
## 60 2024        9      0        ignorado
\end{verbatim}

\hfill\break

Visualizar plot

\begin{Shaded}
\begin{Highlighting}[]
\FunctionTok{ggplot}\NormalTok{(dados\_appended\_melted, }\FunctionTok{aes}\NormalTok{(}\AttributeTok{x=}\NormalTok{year, }\AttributeTok{y=}\NormalTok{value, }\AttributeTok{group=}\NormalTok{raca\_cor\_factor)) }\SpecialCharTok{+}
  \FunctionTok{geom\_line}\NormalTok{(}\FunctionTok{aes}\NormalTok{(}\AttributeTok{linetype=}\NormalTok{raca\_cor\_factor))}\SpecialCharTok{+}
  \FunctionTok{geom\_point}\NormalTok{()}
\end{Highlighting}
\end{Shaded}

\pandocbounded{\includegraphics[keepaspectratio]{bookdown-demo_files/figure-latex/unnamed-chunk-82-1.pdf}}

\begin{Shaded}
\begin{Highlighting}[]
\CommentTok{\# wait to change the variable names:}

\CommentTok{\# names(dados\_appended\_melted)[names(dados\_appended\_melted) == \textquotesingle{}variable\textquotesingle{}] \textless{}{-} \textquotesingle{}RACACOR\textquotesingle{}}
\CommentTok{\# names(dados\_appended\_melted)[names(dados\_appended\_melted) == \textquotesingle{}value\textquotesingle{}] \textless{}{-} \textquotesingle{}n\textquotesingle{}}
\end{Highlighting}
\end{Shaded}

\hfill\break

\section{Mortalidade Anual Total por Neuroblastoma}\label{mortalidade-anual-total-por-neuroblastoma}

\hfill\break

\begin{Shaded}
\begin{Highlighting}[]
\NormalTok{dados\_nb }\OtherTok{\textless{}{-}} \FunctionTok{readRDS}\NormalTok{(}\StringTok{"./data/dados\_nb.rds"}\NormalTok{)}

\CommentTok{\# Checar dados\_nb}
\FunctionTok{head}\NormalTok{(dados\_nb)}
\end{Highlighting}
\end{Shaded}

\begin{verbatim}
##               CONTADOR RACACOR  DTOBITO CAUSABAS   DTNASC IDADE        YMD year
## 876...17551        876       4 20012016     C749 10122010   405 2016-01-20 2016
## 5811...56566    562030       1 14052018     C749 22042018   222 2018-05-14 2018
## 8831...77296    613331       1 11062019     C749 31011964   455 2019-06-11 2019
## 14106...82571   987213       4 01042019     C749 12021950   469 2019-04-01 2019
## 900...87692      78999       1 08032020     C749 02042014   405 2020-03-08 2020
## 8259...95051    470277       1 15022020     C749 21052005   414 2020-02-15 2020
\end{verbatim}

\begin{Shaded}
\begin{Highlighting}[]
\NormalTok{dados\_agregados\_nb\_year }\OtherTok{\textless{}{-}}\NormalTok{ dados\_nb }\SpecialCharTok{\%\textgreater{}\%}
  \FunctionTok{group\_by}\NormalTok{(year) }\SpecialCharTok{\%\textgreater{}\%}
  \FunctionTok{summarise}\NormalTok{(}\AttributeTok{n =} \FunctionTok{n}\NormalTok{())}

\CommentTok{\# Checar dados\_agregados\_nb\_year}
\NormalTok{dados\_agregados\_nb\_year}
\end{Highlighting}
\end{Shaded}

\begin{verbatim}
## # A tibble: 10 x 2
##    year      n
##    <chr> <int>
##  1 2015    310
##  2 2016    261
##  3 2017    306
##  4 2018    322
##  5 2019    313
##  6 2020    201
##  7 2021    243
##  8 2022    346
##  9 2023    329
## 10 2024    352
\end{verbatim}

\begin{Shaded}
\begin{Highlighting}[]
\FunctionTok{ggplot}\NormalTok{(}\AttributeTok{data=}\NormalTok{dados\_agregados\_nb\_year, }\FunctionTok{aes}\NormalTok{(}\AttributeTok{x=}\NormalTok{year, }\AttributeTok{y=}\NormalTok{n, }\AttributeTok{group=}\DecValTok{1}\NormalTok{)) }\SpecialCharTok{+}
  \FunctionTok{geom\_line}\NormalTok{()}\SpecialCharTok{+}
  \FunctionTok{geom\_point}\NormalTok{()}
\end{Highlighting}
\end{Shaded}

\pandocbounded{\includegraphics[keepaspectratio]{bookdown-demo_files/figure-latex/unnamed-chunk-83-1.pdf}}

\hfill\break

\begin{Shaded}
\begin{Highlighting}[]
\NormalTok{dados\_nb\_race\_year\_selcd }\OtherTok{\textless{}{-}}\NormalTok{ dados\_nb }\SpecialCharTok{\%\textgreater{}\%}
  \FunctionTok{select}\NormalTok{(RACACOR, year)}

\CommentTok{\# Checar dados\_nb\_race\_year\_selcd}
\FunctionTok{head}\NormalTok{(dados\_nb\_race\_year\_selcd)}
\end{Highlighting}
\end{Shaded}

\begin{verbatim}
##               RACACOR year
## 876...17551         4 2016
## 5811...56566        1 2018
## 8831...77296        1 2019
## 14106...82571       4 2019
## 900...87692         1 2020
## 8259...95051        1 2020
\end{verbatim}

\begin{Shaded}
\begin{Highlighting}[]
\CommentTok{\# dados\_nb\_race\_2013\_selcd \textless{}{-} dados\_estados\_nb\_2013 \%\textgreater{}\%}
\CommentTok{\#   select(RACACOR, CAUSABAS)}
\end{Highlighting}
\end{Shaded}

\hfill\break

use cast() to convert data frame from long to wide format

\begin{Shaded}
\begin{Highlighting}[]
\FunctionTok{library}\NormalTok{(reshape2)}
\end{Highlighting}
\end{Shaded}

\begin{verbatim}
## 
## Attaching package: 'reshape2'
\end{verbatim}

\begin{verbatim}
## The following object is masked from 'package:tidyr':
## 
##     smiths
\end{verbatim}

\begin{verbatim}
## The following objects are masked from 'package:maditr':
## 
##     dcast, melt
\end{verbatim}

\begin{Shaded}
\begin{Highlighting}[]
\NormalTok{wide\_dados\_nb\_appended }\OtherTok{\textless{}{-}} \FunctionTok{dcast}\NormalTok{(dados\_nb\_race\_year\_selcd, year }\SpecialCharTok{\textasciitilde{}}\NormalTok{ RACACOR)}
\end{Highlighting}
\end{Shaded}

\begin{verbatim}
## Using year as value column: use value.var to override.
\end{verbatim}

\begin{verbatim}
## Aggregation function missing: defaulting to length
\end{verbatim}

\hfill\break

Melt the data frame

\begin{Shaded}
\begin{Highlighting}[]
\NormalTok{dados\_nb\_appended\_melted }\OtherTok{\textless{}{-}} \FunctionTok{melt}\NormalTok{(wide\_dados\_nb\_appended, }\AttributeTok{id.vars =} \StringTok{"year"}\NormalTok{)}
\end{Highlighting}
\end{Shaded}

\hfill\break

Obter variável raça/cor como palavras descritivas ao invés de números.

\begin{Shaded}
\begin{Highlighting}[]
\NormalTok{dados\_nb\_appended\_melted }\OtherTok{\textless{}{-}}\NormalTok{ dados\_nb\_appended\_melted }\SpecialCharTok{\%\textgreater{}\%}
\NormalTok{  dplyr}\SpecialCharTok{::}\FunctionTok{mutate}\NormalTok{(}\AttributeTok{raca\_cor\_factor =} \FunctionTok{ifelse}\NormalTok{(dados\_nb\_appended\_melted}\SpecialCharTok{$}\NormalTok{variable }\SpecialCharTok{==} \StringTok{"1"}\NormalTok{, }\StringTok{"branca"}\NormalTok{, }
                                         \FunctionTok{ifelse}\NormalTok{(dados\_nb\_appended\_melted}\SpecialCharTok{$}\NormalTok{variable }\SpecialCharTok{==} \StringTok{"2"}\NormalTok{, }\StringTok{"preta"}\NormalTok{,}
                                                \FunctionTok{ifelse}\NormalTok{(dados\_nb\_appended\_melted}\SpecialCharTok{$}\NormalTok{variable }\SpecialCharTok{==} \StringTok{"3"}\NormalTok{, }\StringTok{"amarela"}\NormalTok{,}
                                                       \FunctionTok{ifelse}\NormalTok{(dados\_nb\_appended\_melted}\SpecialCharTok{$}\NormalTok{variable }\SpecialCharTok{==} \StringTok{"4"}\NormalTok{, }\StringTok{"parda"}\NormalTok{, }
                                                              \FunctionTok{ifelse}\NormalTok{(dados\_nb\_appended\_melted}\SpecialCharTok{$}\NormalTok{variable }\SpecialCharTok{==} \StringTok{"5"}\NormalTok{, }\StringTok{"indígena"}\NormalTok{,}
                                                                     \FunctionTok{ifelse}\NormalTok{(dados\_nb\_appended\_melted}\SpecialCharTok{$}\NormalTok{variable }\SpecialCharTok{==} \StringTok{"9"}\NormalTok{, }\StringTok{"ignorado"}\NormalTok{,}\StringTok{"0"}\NormalTok{)))))))}
\end{Highlighting}
\end{Shaded}

\hfill\break

Remove NA in year

\begin{Shaded}
\begin{Highlighting}[]
\NormalTok{dados\_nb\_appended\_melted }\OtherTok{\textless{}{-}}\NormalTok{ dados\_nb\_appended\_melted[ }\FunctionTok{which}\NormalTok{(}
\NormalTok{  dados\_nb\_appended\_melted}\SpecialCharTok{$}\NormalTok{year }\SpecialCharTok{!=} \StringTok{"NA"}\NormalTok{), ]}
\end{Highlighting}
\end{Shaded}

\hfill\break

Remove 0 in raça/cor

\begin{Shaded}
\begin{Highlighting}[]
\NormalTok{dados\_nb\_appended\_melted }\OtherTok{\textless{}{-}}\NormalTok{ dados\_nb\_appended\_melted[ }\FunctionTok{which}\NormalTok{(}
\NormalTok{  dados\_nb\_appended\_melted}\SpecialCharTok{$}\NormalTok{raca\_cor\_factor }\SpecialCharTok{!=} \StringTok{"0"}\NormalTok{), ]}
\end{Highlighting}
\end{Shaded}

\hfill\break

Visualizar plot

\begin{Shaded}
\begin{Highlighting}[]
\FunctionTok{ggplot}\NormalTok{(dados\_nb\_appended\_melted, }\FunctionTok{aes}\NormalTok{(}\AttributeTok{x=}\NormalTok{year, }\AttributeTok{y=}\NormalTok{value, }\AttributeTok{group=}\NormalTok{raca\_cor\_factor, }\AttributeTok{color=}\NormalTok{raca\_cor\_factor)) }\SpecialCharTok{+}
  \FunctionTok{geom\_line}\NormalTok{(}\FunctionTok{aes}\NormalTok{(}\AttributeTok{linetype=}\NormalTok{raca\_cor\_factor))}\SpecialCharTok{+}
  \FunctionTok{geom\_point}\NormalTok{()}
\end{Highlighting}
\end{Shaded}

\pandocbounded{\includegraphics[keepaspectratio]{bookdown-demo_files/figure-latex/unnamed-chunk-89-1.pdf}}

\hfill\break

Visualizar Plot Disparity

\begin{Shaded}
\begin{Highlighting}[]
\FunctionTok{library}\NormalTok{(}\StringTok{\textquotesingle{}readxl\textquotesingle{}}\NormalTok{)}
\NormalTok{Disparity\_coeficient\_Brazil\_All }\OtherTok{\textless{}{-}} \FunctionTok{read\_excel}\NormalTok{(}\StringTok{"./data/Disparity\_coeficient\_Brazil\_All.xlsx"}\NormalTok{)}
\end{Highlighting}
\end{Shaded}

\begin{Shaded}
\begin{Highlighting}[]
\FunctionTok{ggplot}\NormalTok{(Disparity\_coeficient\_Brazil\_All, }\FunctionTok{aes}\NormalTok{(}\AttributeTok{x=}\NormalTok{year, }\AttributeTok{y=}\NormalTok{value, }\AttributeTok{group=}\NormalTok{racacor, }\AttributeTok{color=}\NormalTok{racacor)) }\SpecialCharTok{+}
  \FunctionTok{geom\_line}\NormalTok{(}\FunctionTok{aes}\NormalTok{(}\AttributeTok{linetype=}\NormalTok{racacor))}\SpecialCharTok{+}
  \FunctionTok{geom\_point}\NormalTok{()}
\end{Highlighting}
\end{Shaded}

\pandocbounded{\includegraphics[keepaspectratio]{bookdown-demo_files/figure-latex/unnamed-chunk-91-1.pdf}}

\hfill\break

Visualizar Plot Disparity

\begin{Shaded}
\begin{Highlighting}[]
\FunctionTok{library}\NormalTok{(}\StringTok{\textquotesingle{}readxl\textquotesingle{}}\NormalTok{)}
\NormalTok{Disparity\_coeficient\_Sum\_PardosBlacks\_Brazil\_All }\OtherTok{\textless{}{-}} \FunctionTok{read\_excel}\NormalTok{(}\StringTok{"./data/Disparity\_coeficient\_Sum\_PardosBlacks\_Brazil\_All.xlsx"}\NormalTok{)}
\end{Highlighting}
\end{Shaded}

\begin{Shaded}
\begin{Highlighting}[]
\FunctionTok{ggplot}\NormalTok{(Disparity\_coeficient\_Sum\_PardosBlacks\_Brazil\_All, }\FunctionTok{aes}\NormalTok{(}\AttributeTok{x=}\NormalTok{year, }\AttributeTok{y=}\NormalTok{value, }\AttributeTok{group=}\NormalTok{racacor, }\AttributeTok{color=}\NormalTok{racacor)) }\SpecialCharTok{+}
  \FunctionTok{geom\_line}\NormalTok{(}\FunctionTok{aes}\NormalTok{(}\AttributeTok{linetype=}\NormalTok{racacor))}\SpecialCharTok{+}
  \FunctionTok{geom\_point}\NormalTok{()}
\end{Highlighting}
\end{Shaded}

\pandocbounded{\includegraphics[keepaspectratio]{bookdown-demo_files/figure-latex/unnamed-chunk-93-1.pdf}}

\hfill\break

Visualizar Plot Disparity IBGE

\begin{Shaded}
\begin{Highlighting}[]
\FunctionTok{library}\NormalTok{(}\StringTok{\textquotesingle{}readxl\textquotesingle{}}\NormalTok{)}
\NormalTok{Disparity\_ibge\_plot }\OtherTok{\textless{}{-}} \FunctionTok{read\_excel}\NormalTok{(}\StringTok{"./data/Disparity\_ibge\_plot.xlsx"}\NormalTok{)}
\end{Highlighting}
\end{Shaded}

\begin{Shaded}
\begin{Highlighting}[]
\FunctionTok{ggplot}\NormalTok{(Disparity\_ibge\_plot, }\FunctionTok{aes}\NormalTok{(}\AttributeTok{x=}\NormalTok{year, }\AttributeTok{y=}\NormalTok{value, }\AttributeTok{group=}\NormalTok{racacor, }\AttributeTok{color=}\NormalTok{racacor)) }\SpecialCharTok{+}
  \FunctionTok{geom\_line}\NormalTok{(}\FunctionTok{aes}\NormalTok{(}\AttributeTok{linetype=}\NormalTok{racacor))}\SpecialCharTok{+}
  \FunctionTok{geom\_point}\NormalTok{()}
\end{Highlighting}
\end{Shaded}

\begin{verbatim}
## Warning: Removed 10 rows containing missing values or values outside the scale range
## (`geom_line()`).
\end{verbatim}

\begin{verbatim}
## Warning: Removed 10 rows containing missing values or values outside the scale range
## (`geom_point()`).
\end{verbatim}

\pandocbounded{\includegraphics[keepaspectratio]{bookdown-demo_files/figure-latex/unnamed-chunk-95-1.pdf}}

\hfill\break

Visualizar Plot Disparity IBGE

\begin{Shaded}
\begin{Highlighting}[]
\FunctionTok{library}\NormalTok{(}\StringTok{\textquotesingle{}readxl\textquotesingle{}}\NormalTok{)}
\NormalTok{Disparity\_ibge\_negros\_plot }\OtherTok{\textless{}{-}} \FunctionTok{read\_excel}\NormalTok{(}\StringTok{"./data/Disparity\_ibge\_negros\_plot.xlsx"}\NormalTok{)}
\end{Highlighting}
\end{Shaded}

\begin{Shaded}
\begin{Highlighting}[]
\FunctionTok{ggplot}\NormalTok{(Disparity\_ibge\_negros\_plot, }\FunctionTok{aes}\NormalTok{(}\AttributeTok{x=}\NormalTok{year, }\AttributeTok{y=}\NormalTok{value, }\AttributeTok{group=}\NormalTok{racacor, }\AttributeTok{color=}\NormalTok{racacor)) }\SpecialCharTok{+}
  \FunctionTok{geom\_line}\NormalTok{(}\FunctionTok{aes}\NormalTok{(}\AttributeTok{linetype=}\NormalTok{racacor))}\SpecialCharTok{+}
  \FunctionTok{geom\_point}\NormalTok{()}
\end{Highlighting}
\end{Shaded}

\pandocbounded{\includegraphics[keepaspectratio]{bookdown-demo_files/figure-latex/unnamed-chunk-97-1.pdf}}

\hfill\break

Stacked Graph

\begin{Shaded}
\begin{Highlighting}[]
\FunctionTok{ggplot}\NormalTok{(dados\_nb\_appended\_melted, }\FunctionTok{aes}\NormalTok{(}\AttributeTok{fill=}\NormalTok{raca\_cor\_factor, }\AttributeTok{y=}\NormalTok{value, }\AttributeTok{x=}\NormalTok{year)) }\SpecialCharTok{+} 
    \FunctionTok{geom\_bar}\NormalTok{(}\AttributeTok{position=}\StringTok{"stack"}\NormalTok{, }\AttributeTok{stat=}\StringTok{"identity"}\NormalTok{)}
\end{Highlighting}
\end{Shaded}

\pandocbounded{\includegraphics[keepaspectratio]{bookdown-demo_files/figure-latex/unnamed-chunk-98-1.pdf}}

\hfill\break

Stacked + Graph

\begin{Shaded}
\begin{Highlighting}[]
\FunctionTok{ggplot}\NormalTok{(dados\_nb\_appended\_melted, }\FunctionTok{aes}\NormalTok{(}\AttributeTok{fill=}\NormalTok{raca\_cor\_factor, }\AttributeTok{y=}\NormalTok{value, }\AttributeTok{x=}\NormalTok{year)) }\SpecialCharTok{+} 
    \FunctionTok{geom\_bar}\NormalTok{(}\AttributeTok{position=}\StringTok{"fill"}\NormalTok{, }\AttributeTok{stat=}\StringTok{"identity"}\NormalTok{)}
\end{Highlighting}
\end{Shaded}

\pandocbounded{\includegraphics[keepaspectratio]{bookdown-demo_files/figure-latex/unnamed-chunk-99-1.pdf}}

\chapter{Final Words}\label{final-words}

We have finished a nice book.

  \bibliography{book.bib,packages.bib}

\end{document}
